% =====================================================================
%  PAPER D  --  merged submission for NDSS (Network and Distributed
%  System Security Symposium), combining the cross-layer benchmark
%  (formerly Paper C) and the composition theorem (formerly Paper A).
%
%  NDSS FORMAT COMPLIANCE
%  ----------------------
%  NDSS requires: US Letter (not A4), two-column, columns no more than
%  9.25 in high and 3.5 in wide, Times font at 10 pt or larger with
%  11 pt or larger leading, and at most 13 pages excluding the Ethics
%  Considerations section, references, and appendices.
%
%  The class line below (IEEEtran, conference, letterpaper, 10pt, with
%  Times via mathptmx) reproduces exactly that geometry: \columnwidth
%  = 3.5 in and \textheight = 9.25 in (verified at compile time by the
%  \typeout block after \begin{document}).
%
%  FOR FINAL SUBMISSION: download the official template from
%  ndss-symposium.org/ndss2026/submissions/templates/ and replace the
%  two lines marked "NDSS SWAP" below with \documentclass{ndss}.
%  No change to the body is required.
% =====================================================================
\documentclass[conference,letterpaper,10pt]{IEEEtran}   % NDSS SWAP (1/2)
\IEEEoverridecommandlockouts
\usepackage{mathptmx}                                    % NDSS SWAP (2/2): Times
\usepackage{cite}
\usepackage{amsmath,amssymb,amsfonts}
\usepackage{amsthm}
\usepackage{graphicx}
\usepackage{booktabs}
\usepackage{tabularx}
\usepackage{subcaption}
\usepackage{xcolor}
\usepackage{tikz}
\usepackage{pgfplots}
\pgfplotsset{compat=1.17}
\usetikzlibrary{arrows.meta,positioning,fit,backgrounds,shapes.geometric,calc}
\usepackage[hidelinks]{hyperref}

% NDSS caps column height at 9.25in; IEEEtran's default is 9.30in, so we
% trim it to be strictly within spec (9.25in = 668.5pt at 72.27pt/in).
\setlength{\textheight}{668.4pt}

% ---- anonymisation switch: NDSS review is double-blind ----------------
% \anontrue  = anonymised submission version (what NDSS requires)
% \anonfalse = named version for internal / supervisor review
\newif\ifanon
\anontrue

% ---- shared style tokens --------------------------------------------
\definecolor{netcol}{RGB}{31,78,121}
\definecolor{autocol}{RGB}{18,120,90}
\definecolor{bridgecol}{RGB}{176,96,0}
\definecolor{accent}{RGB}{140,20,20}
\newcommand{\layernet}{\textcolor{netcol}{\textbf{network}}}
\newcommand{\layerauto}{\textcolor{autocol}{\textbf{autonomy}}}
\newcommand{\dataset}[1]{\textsc{#1}}
\newcommand{\code}[1]{\texttt{#1}}
\newcommand{\Real}{\mathbb{R}}
\newcommand{\Norm}[1]{\left\lVert #1\right\rVert}
\newcommand{\MCR}{\mathrm{MCR}}

\newtheorem{theorem}{Theorem}
\newtheorem{lemma}{Lemma}
\newtheorem{definition}{Definition}
\newtheorem{proposition}{Proposition}

\begin{document}

% Title states ONE thesis: the end-to-end framework. The corpus is the
% framework's instrument, not a co-equal deliverable, and the attack class the
% theorem covers (time-delay forwarding) is named rather than left as a vague
% "under attack".
\title{From Wire to Flight: An End-to-End Framework Composing\\
Network Intrusion Detection with Certified Mission Safety\\
for Autonomous UAVs}

\ifanon
  \author{\IEEEauthorblockN{Paper \#XXX}
  \IEEEauthorblockA{Submitted for double-blind review}}
\else
  \author{
  \IEEEauthorblockN{Roger Nick Anaedevha\IEEEauthorrefmark{1}, \textit{IEEE Member},
  Keiwan Soltani\IEEEauthorrefmark{2}, and Federico Cor\`o\IEEEauthorrefmark{3}}
  \IEEEauthorblockA{\IEEEauthorrefmark{1}Institute of Cyber Intelligent Systems,
  National Research Nuclear University MEPhI, Moscow, Russia}
  \IEEEauthorblockA{\IEEEauthorrefmark{2}Computer Science Department,
  Missouri University of Science and Technology, Rolla, MO, USA}
  \IEEEauthorblockA{\IEEEauthorrefmark{3}Tenure-Track Assistant Professor (RTT),
  University of Padova, Padua, Italy}}
\fi

\maketitle

% compile-time proof of NDSS geometry conformance (appears in the log)
\typeout{NDSS GEOMETRY CHECK: columnwidth=\the\columnwidth\space
(target 3.5in=252.94pt), textheight=\the\textheight\space
(target 9.25in=668.62pt)}

\begin{abstract}
Unmanned aerial vehicle (UAV) security is presently studied as two disjoint
problems. Network intrusion detection bounds what reaches the aircraft over the
swarm mesh, the command link, and the on-board bus, but makes no statement about
the flight; certified robustness bounds a navigation model's response to a
perturbation of its input, but treats the origin of that perturbation as
exogenous. Neither discipline can therefore answer the question that determines
airworthiness: whether an attack that evades the detector still permits the
mission to complete. This paper contributes an end-to-end framework in which
that question is both well posed and provable, and it is the framework---not any
single dataset---that is the contribution. The framework has three components.
A two-layer schema promotes the \emph{cross-layer pairing}---a network attack
window joined to the navigation degradation it induces---to a first-class
object, carrying a mandatory provenance flag recording whether the coupling was
measured on one platform, physically modelled, or aligned in simulation. An
instrumented corpus of six datasets spanning both layers, together with a
harness that exposes a detector's operating point as a swept parameter, supplies
the quantities the framework consumes. On these objects we prove a
\emph{composition theorem}: for any time-delay forwarding attack that evades a
detector at operating point $\theta$, the navigation stack retains a certified
Mission-Completion-Rate $\MCR\ge f(\delta(\theta))$. Two measurements make the
guarantee sharp rather than nominal: the residual attack budget \emph{saturates}
at the network's contact-window slack, at most $4.84$\,s of undetected delay
over $15{,}392$ observed hops, and the delay-to-position interface, calibrated
on real GPS-spoofing and jamming flights from the target autopilot family, is an
order of magnitude tighter than the kinematic worst case. Under that measured
interface the detector's operating point lies inside a nonempty certified
window, and the composed configuration is the only one that certifies a nonzero
mission floor against an evading attacker, dominating both single-layer
baselines at every operating point (Wilcoxon signed-rank, Holm-corrected,
$p<10^{-20}$). Schema, adapters, harness, and certificate engine are released as
an artifact in which every reported number carries a provenance tag.
\end{abstract}

% =====================================================================
\section{Introduction}
\label{sec:intro}
% =====================================================================
A modern autonomous UAV is a stack of trust boundaries. Commands descend from a
ground station over a command-and-control (C2) link; peer aircraft exchange
state over a multi-hop swarm mesh; motor controllers and sensors talk over an
on-board bus; and a navigation pipeline fuses satellite, inertial, and visual
measurements into the state estimate that closes the flight control loop. An
adversary can enter at any of these boundaries, and the security literature has
produced strong, specialised defences at each. What it has \emph{not} produced
is a way to reason across them.

This gap is not academic. Electronic-warfare environments defeat drones not by
breaking cryptography but by degrading the physics: a jammer raises the
jamming-to-signal ratio ($J/S$) until the satellite-navigation solution
collapses and the mission fails~\cite{kassas2022jammer}. A forwarding attack in
the mesh does not exfiltrate data; it delays a correction until the aircraft is
acting on stale state~\cite{datamut2026,lou2019timedelay}. In every case the
\emph{network} event and the \emph{navigation} outcome are two ends of one
causal chain---and yet neither the datasets nor the guarantees represent that
chain.

Consider what a reviewer of a UAV intrusion detector can and cannot check today.
They can verify that the detector flags a delayed forwarding event or a spoofed
navigation message with high recall. They \emph{cannot} verify the claim that
matters operationally---that a detector tuned to a given false-alarm budget
keeps aircraft completing their missions---because no dataset records both the
network event and the flight it degrades. Symmetrically, ask a
certified-robustness researcher whether the UAV is safe and the answer is ``the
navigation model is provably stable for input perturbations up to radius $R$'':
a statement about a model, not a flight, and one that never says where a
radius-$R$ perturbation comes from or how often an attacker can produce it.

We close the gap from both ends, and the two halves are inseparable: the
guarantee needs objects the benchmark defines, and the benchmark needs the
guarantee to show what those objects are for.

\emph{The benchmark half.} We unify six datasets---four established
network-layer corpora and two navigation-layer sources---under a two-layer
schema whose unit of analysis is neither a packet nor a flight but a
\emph{cross-layer pairing}: a network attack window $W_n$ joined to an autonomy
degradation window $W_a$, carrying the detector operating point, the residual
undetected budget, the induced perturbation, and the mission outcome. Because a
cross-layer corpus invites exactly one dangerous mistake---asserting that a
network event caused a flight degradation the two were never observed
together---every pairing carries a mandatory \code{pairing\_basis} flag
recording how the coupling was established.

\emph{The guarantee half.} On those objects we prove a chained statement,
\begin{equation}
\label{eq:chain}
\theta \;\longmapsto\; \Delta(\theta) \;\longmapsto\; \delta(\theta)
\;\longmapsto\; \MCR \ge f\big(\delta(\theta)\big),
\end{equation}
read left to right: a detector operating point $\theta$ bounds the residual
undetected attack $\Delta(\theta)$; the residual attack maps to a bounded
perturbation $\delta(\theta)$ on the navigation input; and a
Lipschitz--Gr\"onwall trajectory-tube argument turns that bounded input into a
guaranteed Mission-Completion-Rate (MCR) floor. Each arrow is individually
grounded in prior work---detector thresholds, worst-case attack analysis,
certified robustness---but the \emph{composition} is new: no prior guarantee
chains a network-security threshold to a navigation-safety floor.

\subsection*{Contributions}
\begin{itemize}
\item \textbf{A two-layer schema and six validated adapters
(\S\ref{sec:schema},~\S\ref{sec:corpus}).} A single representation for any
observation from any of the six sources, with lossless ingestion and a
per-record honesty flag that makes the strength of every cross-layer coupling
explicit and auditable.

\item \textbf{A detector-operating-curve harness (\S\ref{sec:harness}).} We make
the network detector's operating point a swept parameter and measure recall and
false-positive rate as a function of it, producing the network-side premise a
composed guarantee consumes. Across three topologies the curve has a universal
knee-and-floor shape whose floor is the attacker's self-exposure probability.

\item \textbf{An empirically tightened interface (\S\ref{sec:interface}).} We
prove a residual-budget lemma and then \emph{tighten it by measurement}: the
undetected budget saturates at the contact-window slack, verified at
$\le 4.84$\,s over $15{,}392$ hops. We calibrate the delay-to-position rate on
three real spoofing and jamming flights, obtaining a bound an order of magnitude
tighter than the kinematic worst case.

\item \textbf{A composition theorem and a four-certificate architecture
(\S\ref{sec:theorem},~\S\ref{sec:certs}).} We prove~\eqref{eq:chain} and its
soundness, and show precisely why the deterministic core does not subsume
randomized smoothing, PAC-Bayes, or multiplicative-weights regret.

\item \textbf{Instantiation, dominance, and an open artifact
(\S\ref{sec:eval}).} The composed guarantee is the only configuration with a
nonzero certified floor for an evading attacker; under the measured interface
the detector's operating point lies inside the certified window. Every figure
and table in this paper is generated from committed result files, and every
number carries a provenance tag.
\end{itemize}

\emph{Organization.} Appendix~\ref{sec:background} gives the system model,
terminology, and symbol table a reader outside UAV security needs; acronyms are
also defined at first use in the body.
\S\ref{sec:problem} states the threat model and formalises the missing object.
\S\ref{sec:related} positions the work. \S\ref{sec:schema}--\S\ref{sec:harness}
build the benchmark; \S\ref{sec:interface}--\S\ref{sec:certs} build the
guarantee; \S\ref{sec:eval} instantiates and evaluates it.
\S\ref{sec:discussion}--\S\ref{sec:conclusion} discuss, bound, and conclude.

% =====================================================================
\section{Threat Model and Problem Formulation}
\label{sec:problem}
% =====================================================================
\subsection{The two attack surfaces}
We partition the UAV attack surface into two \emph{layers}. The \layernet{}
layer covers everything between nodes: the swarm mesh (multi-hop routing among
aircraft), the C2 link, and the intra-UAV bus. The \layerauto{} layer covers
everything inside the navigation and control stack: the GNSS receiver, sensor
fusion, the flight controller, and the mission outcome. A \emph{sublayer}
refines each into the concrete surfaces a dataset can observe.
Fig.~\ref{fig:surfaces} places the six datasets on the causal chain.

\begin{figure}[t]
\centering
\resizebox{\columnwidth}{!}{%
\begin{tikzpicture}[font=\footnotesize,>={Latex[length=2mm]},
  node distance=2.5mm,
  n/.style={rounded corners=2pt,draw=netcol,fill=netcol!10,align=center,
    minimum height=6.5mm,inner sep=2.5pt,text width=17mm},
  a/.style={rounded corners=2pt,draw=autocol,fill=autocol!10,align=center,
    minimum height=6.5mm,inner sep=2.5pt,text width=17mm}]
\node[n] (mesh) {swarm mesh\\\scriptsize UAVIDS, UAV-CAS, DATAMUt};
\node[n,right=of mesh] (c2) {C2 link\\\scriptsize UAV Attack Dataset};
\node[n,right=of c2] (bus) {intra-bus\\\scriptsize HCRL UAVCAN};
\node[a,below=8mm of c2] (gnss) {GNSS + EKF\\\scriptsize state estimate};
\node[a,right=of gnss] (ctrl) {controller\\+ mission\\\scriptsize UAV-EW-Bench};
\node[a,left=of gnss] (nav) {navigation\\input $u(t)$};
% mesh feeds the navigation input directly (peer state corrections). Drawn to
% the node itself; an earlier version terminated on a coordinate expression
% that left the arrowhead pointing into empty space.
\draw[->,netcol] (mesh) -- (nav);
\draw[->,netcol] (c2) -- (gnss);
\draw[->,netcol] (bus) -- (ctrl.north);
\draw[->,autocol] (nav) -- (gnss);
\draw[->,autocol] (gnss) -- (ctrl);
\node[align=center,font=\scriptsize\itshape] at ($(mesh)!0.5!(bus)+(0,0.9)$)
  {network layer (attacks \emph{on the wire})};
\node[align=center,font=\scriptsize\itshape] at ($(nav)!0.5!(ctrl)-(0,0.95)$)
  {autonomy layer (attacks \emph{on the flight})};
\end{tikzpicture}%
}
\caption{The attack surfaces the six datasets cover, arranged on the causal
chain. Network-layer events (top) propagate downward into the autonomy layer
(bottom) through the navigation input; representing that downward arrow
explicitly, rather than assuming it, is what the benchmark contributes and what
the theorem then exploits.}
\label{fig:surfaces}
\end{figure}

\subsection{Adversary and detector}
We assume an adversary who can (i) inject, drop, delay, or forge traffic on one
or more network sublayers, and (ii) manipulate the RF environment of the GNSS
receiver. We assume a network-layer \emph{detector} $D_\theta$ operating at a
tunable point $\theta$ that trades detection recall against false-positive rate:
for the multi-hop forwarding detector we build on, $\theta$ is a per-hop
deviation tolerance in seconds, and a hop is flagged if its forwarding delay
exceeds $\theta$ or its next hop deviates from a scheduling oracle. Lowering
$\theta$ catches more but raises false positives; raising it does the reverse.
An \emph{evading} attacker chooses its behaviour to stay just under $\theta$ and
is therefore invisible by construction. We assume the adversary knows $\theta$;
this is the strongest reasonable assumption and the one our worst-case analysis
requires.

\subsection{Mission completion: a spatial and a temporal predicate}
\label{sec:mcr}
We adopt \emph{Mission Completion Rate} (MCR) as the operational metric, and
take $\MCR\ge0.90$ as a reference consistent with the airworthiness-security
posture of DO-326A/ED-202A~\cite{do326a}; emerging AI-in-aviation assurance
guidance~\cite{ed324} motivates treating a certified MCR floor as a first-class
safety property, which is what \S\ref{sec:theorem} delivers.

A purely \emph{spatial} definition of completion---the trajectory never leaves
the safe corridor---is, however, inadequate for precisely the attack class this
paper targets. A time-delay adversary can leave the aircraft geometrically on
route while making it arrive late; a delivery, inspection, or search mission
with a time-critical objective has then failed even though no corridor was ever
violated. We therefore define completion as the \emph{conjunction} of a spatial
and a temporal predicate. Writing $x(t)$ for the realised trajectory,
$x_{\mathrm{nom}}(t)$ for the nominal one, $m$ for the corridor margin,
$t_{\mathrm{arr}}$ for the arrival time at the objective, $T$ for the nominal
mission duration, and $\kappa$ for the mission's schedule tolerance
(a fraction of $T$),
\begin{equation}
\label{eq:mcr}
\MCR \;=\; \Pr\Big[\underbrace{\sup_{t\in[0,T]}\Norm{x(t)-x_{\mathrm{nom}}(t)}\le m}_{\text{spatial}}
\;\wedge\;
\underbrace{t_{\mathrm{arr}}\le(1+\kappa)T}_{\text{temporal}}\Big].
\end{equation}
Where a quantity is computed under the spatial predicate alone---as the
navigation anchor of \S\ref{sec:eval} is, because its source records geometric
completion only---we label it \emph{Spatial MCR} and say so explicitly.

The temporal predicate is not a separate analysis: the same residual budget
$\Delta(\theta)$ that bounds the spatial perturbation also bounds the schedule
slip, since an evading attacker's total contribution to arrival delay is exactly
its accumulated undetected forwarding delay. The temporal predicate therefore
holds whenever
\begin{equation}
\label{eq:temporal_condition}
\Delta(\theta)\;\le\;\kappa\,T ,
\end{equation}
which by Lemma~\ref{lem:residual} is implied by $H\min(\theta,s)\le\kappa T$.
This adds a second edge to the certified window of \S\ref{sec:eval}, and it is
worth recording which edge binds. At the operating point we certify
($\theta=0.25$\,s, $H=2$) the induced slip is at most $0.5$\,s, so for any
mission longer than $5$\,s at a $10\%$ schedule tolerance the temporal predicate
is satisfied with room to spare and the spatial predicate is the binding one.
The ordering reverses for a loosely tuned detector: at $\theta$ near the
contact-window slack $s=5$\,s the slip reaches $H s = 10$\,s, which exceeds a
$10\%$ tolerance on any mission shorter than $100$\,s. A short, time-critical
mission is thus governed by \eqref{eq:temporal_condition} and a long, tightly
corridored one by the spatial bound---a distinction a spatial-only metric
cannot express.

\subsection{The missing object}
\label{sec:formulation}
Let a network attack occupy a time window $W_n=[t_0,t_1]$ on some network
sublayer, and let the induced navigation degradation occupy a window $W_a$ on
the autonomy layer. Existing datasets provide either $W_n$ (with packet labels)
or $W_a$ (with flight outcomes) but never the \emph{pair} together with the
mechanism linking them. We formalise the missing object as a tuple
\begin{equation}
\label{eq:pairing}
P = \big(W_n,\; W_a,\; \theta,\; \Delta(\theta),\; \delta,\; \MCR\big),
\end{equation}
where $\theta$ is the detector operating point, $\Delta(\theta)$ the worst-case
residual effect of an attacker that evades $D_\theta$, $\delta$ the resulting
perturbation budget on the navigation input, and MCR the measured or certified
outcome. Representing $P$---not just $W_n$ or $W_a$---is the benchmark's
contribution; bounding the last component of $P$ from the first is the
theorem's. Crucially, the honesty of $P$ depends on \emph{how} $W_n$ and $W_a$
were linked, which we record explicitly (\S\ref{sec:schema}).

\begin{definition}[Composed guarantee]
\label{def:composed}
Given a detector $D_\theta$ and navigation dynamics \eqref{eq:ode}, a
\emph{composed guarantee} is a function $f$ such that for every attack $a$ with
$D_\theta(a)=\text{``clean''}$, the mission completes with rate
$\MCR(a)\ge f(\delta(\theta))$, where $\delta(\theta)$ bounds the
navigation-input perturbation any such $a$ can induce.
\end{definition}

% =====================================================================
\section{Related Work}
\label{sec:related}
% =====================================================================
Prior work divides along the two halves of this paper---corpora on the data
side, bounds on the guarantee side---and the same gap recurs on both: each
community works within one layer and leaves the interface between them to
assumption.

\paragraph{UAV corpora, and the discipline of building them}
Taxonomies of AI-enhanced UAV intrusion detection~\cite{islam2025uavids} report
a proliferation of datasets alongside a persistent weakness: general-purpose
network corpora lack UAV protocol semantics, while UAV-specific corpora each
cover a single sublayer. The consequence is measurable---models trained on
generic benchmarks misclassify a large fraction of GPS-spoofing events on real
autopilot logs~\cite{islam2025uavids}---and it is direct evidence that the
network and navigation layers do not form one undifferentiated feature space.
FANET surveys~\cite{tsao2022survey,ceviz2024survey} catalogue routing attacks
but stop at network effect. The individual corpora are authoritative within
their sublayers: \dataset{UAVIDS-2025}~\cite{uavids2025,zeng2025uavids}
contributes $122{,}171$ NS-3 mesh flows, the \dataset{UAV Attack
Dataset}~\cite{whelan2020uav,whelan2020novelty} real PX4 flights under GPS
spoofing, jamming, and MAVLink ping-DoS, and \dataset{HCRL
UAVCAN}~\cite{hcrl_uavcan} bus-level flooding, fuzzing, and replay. None reaches
the flight.

How such a corpus is built matters as much as how large it is. Audits of
CIC-IDS2017 exposed labelling errors, flawed flow construction, and unrealistic
temporal separation, each of which inflated reported
performance~\cite{engelen2021troubleshooting,sharafaldin2018cicids}. Spanning
two layers introduces a further failure mode: asserting a causal link between a
network event and a flight degradation that were never observed together. The
mandatory provenance flag of \S\ref{sec:schema} exists to bound precisely that
risk. Among released datasets the closest in ambition is
\dataset{UAV-CAS}~\cite{uavcas2026}, an AERPAW-calibrated swarm corpus and the
first to include collaborative multi-attack compositions; we ingest it as one of
our six sources. It nonetheless remains a network-layer benchmark---its authors
list GNSS spoofing and jamming as future work---so where it asks how well a
collaborative mesh attack can be detected, we ask what that attack, detected or
missed, does to the aircraft.

\paragraph{Certified robustness, and its lift to control}
Randomized smoothing~\cite{cohen2019smoothing,salman2019provably} yields a
probabilistic $\ell_2$ radius that scales to large models, whereas
Lipschitz-based certification~\cite{pauli2022training,zhang2022rethinking}
yields deterministic bounds known to be loose globally; neural
ODEs~\cite{chen2018neuralode} admit the Gr\"onwall-type
argument~\cite{yan2020robustnode} we adopt. All of these bound a \emph{model's}
output. Lifting them to sequential decisions, CROP certifies a cumulative-reward
lower bound under state perturbation~\cite{wu2022crop}, CARRL certifies robust
actions~\cite{everett2021carrl}, reachability tools bound trajectories of
network-controlled systems~\cite{ivanov2019verisig,tran2020nnv}, and
control-theoretic work establishes safety under sensor false-data
injection~\cite{li2022lqg}. CROP is the nearest analogue to a certified mission
floor, but it certifies a policy against a \emph{given} perturbation ball and is
silent on where that ball comes from.

\paragraph{Estimation, assurance, and delay}
A parallel control-theoretic literature asks how much sensor corruption an
estimator can tolerate: exact recovery holds below a corruption
threshold~\cite{fawzi2014secure}, satisfiability-modulo-theory solvers make the
combinatorial search practical~\cite{shoukry2017smt}, and recent surveys
catalogue the field~\cite{ding2021survey}. Such results bound estimation error
given \emph{which} sensors are compromised, rather than modelling a detector
whose threshold determines what an attacker may inject at all---the premise
Lemma~\ref{lem:residual} supplies. Simplex and its
descendants~\cite{sha2001simplex,mehmood2022blackbox,cofer2020rta} guarantee
safety by switching from an unverified controller to a verified baseline when a
monitor fires; our certified window is precisely the region in which the
high-performance policy may run, and we contribute the quantitative boundary
that such architectures otherwise hand-tune. The destabilising effect of
injected delay is itself well studied, with detectors built on state-aware
inference and divergence tests~\cite{lou2019timedelay,etd2023uav,datamut2026},
but that literature establishes whether a delay attack is dangerous or
detectable and stops there.

\paragraph{Composing guarantees}
The closest prior art composes certificates \emph{within} a single ML pipeline.
Sheikholeslami et al.~\cite{sheikholeslami2021provably} jointly certify a
classifier and a detection class, so that an input is provably either flagged or
correctly classified---structurally our evade-or-bound dichotomy, but confined
to one layer. Yang et al.~\cite{yang2022sensing} decompose end-to-end certified
robustness across a two-stage sensing$\to$reasoning pipeline, the strongest
existing template for chained certification, though it too remains within
perception. Assume-guarantee and contract-based CPS design~\cite{pacti2024}
supplies the compositional vocabulary but no ML-robustness certificates. No
prior result chains a network-intrusion-detection guarantee to a navigation
certificate; that is the gap this paper closes, and closing it requires the
corpus and the theorem together.

% =====================================================================
\section{The Two-Layer Schema}
\label{sec:schema}
% =====================================================================
The schema has three record kinds, designed so the cross-layer pairing is a
first-class object rather than an implicit join.

\textbf{Event} is one atomic observation---a mesh flow, a bus frame, a routing
hop, a telemetry sample, or a whole flight---mapped to a small common core
(identity, layer/sublayer, relative time, endpoints, a unified attack label)
plus a \code{metrics} block of the few cross-source numeric quantities that
matter (residual forwarding delay, position error, $J/S$, mission-completed) and
a \code{native} block preserving every remaining source column, so ingestion is
lossless.

\textbf{AttackWindow} is a labelled attack interval within one scenario,
carrying the unified attack class and class-appropriate \code{intensity}
parameters. It is the unit of cross-layer alignment.

\textbf{CrossLayerPairing} realises the tuple $P$ of Eq.~\eqref{eq:pairing}: it
joins a network window to an autonomy window and records $\theta$, the residual
budget $\Delta(\theta)$, the perturbation $\delta$ with a \emph{named,
versioned} mapping, the certificate hook, and the empirical MCR. Its
\code{pairing\_basis} field is the schema's conscience:
\code{measured\_same\_platform} (both windows recorded on one testbed---the
strongest evidence), \code{physics\_model} (coupled through a stated physical
model), or \code{synthetic\_alignment} (aligned by construction). A reviewer can
filter the benchmark by evidence strength; a claim can never silently upgrade
its own provenance.

Three design choices let six heterogeneous sources coexist without information
loss. Time is always \emph{relative} to each source's own capture start, because
the sources share no global clock; cross-layer alignment happens at the window
level, so a mismatch between a $50$\,Hz telemetry stream and a per-flow network
record never forces a lossy resample. Attack labels are \emph{unified} into a
common taxonomy while the verbatim source label is retained in
\code{label\_native}, so a model can train on the unified class and an auditor
can still recover what the source called it. And every unmapped column is
preserved, which makes the ingestion provably lossless rather than merely
convenient.

Cross-layer alignment therefore needs a stated rule and a stated error budget,
since this is the step at which a benchmark could silently manufacture---or
destroy---the effect it reports. We anchor attack onsets under an affine map
between the two relative clocks and bound the residual misalignment by
$e_{\mathrm{align}}\le0.53$\,s for every released pairing; crucially, the
quantity the certificate consumes, $\Delta(\theta)$, is computed entirely within
the network source, so misalignment can attach a pairing to the wrong autonomy
window but cannot inflate or deflate $\Delta(\theta)$ itself.
Appendix~\ref{app:alignment} gives the map, the three error terms, and the
rejection rule.

\begin{figure}[t]
\centering
\resizebox{\columnwidth}{!}{%
\begin{tikzpicture}[
  font=\footnotesize,
  box/.style={rounded corners=2pt,draw,align=center,minimum height=7mm,inner sep=3pt},
  net/.style={box,fill=netcol!12,draw=netcol},
  aut/.style={box,fill=autocol!12,draw=autocol},
  br/.style={box,fill=bridgecol!15,draw=bridgecol,thick},
  ev/.style={box,fill=black!4,minimum width=15mm},
  >={Latex[length=2mm]}]
\node[ev] (mesh) {mesh\\flow};
\node[ev,right=3mm of mesh] (c2) {C2\\frame};
\node[ev,right=3mm of c2] (bus) {bus\\frame};
\node[ev,right=3mm of bus] (hop) {route\\hop};
\node[ev,right=3mm of hop] (tel) {telem.\\sample};
\node[ev,right=3mm of tel] (fl) {flight};
\node[left=2mm of mesh] (evlab) {\textbf{Event}};
\node[net,below=7mm of c2,minimum width=30mm] (wn) {network AttackWindow $W_n$};
\node[aut,below=7mm of tel,minimum width=30mm] (wa) {autonomy AttackWindow $W_a$};
\node[br,below=8mm of $(wn)!0.5!(wa)$,minimum width=62mm] (pair)
  {\textbf{CrossLayerPairing} $P=(W_n,W_a,\theta,\Delta(\theta),\delta,\mathrm{MCR})$\\
   \scriptsize \code{pairing\_basis} $\in$ \{measured, physics-model, synthetic\}};
\draw[->] (mesh) -- (wn); \draw[->] (c2) -- (wn); \draw[->] (bus) -- (wn); \draw[->] (hop) -- (wn);
\draw[->] (tel) -- (wa); \draw[->] (fl) -- (wa);
\draw[->,bridgecol,thick] (wn) -- (pair);
\draw[->,bridgecol,thick] (wa) -- (pair);
\end{tikzpicture}%
}
\caption{The schema in one picture. Heterogeneous events from six datasets fold
into typed network and autonomy \emph{AttackWindows}; a \emph{CrossLayerPairing}
joins one of each and records the detector operating point, the residual budget,
the induced perturbation, and the mission outcome, tagged by how the coupling
was established.}
\label{fig:schema}
\end{figure}

% =====================================================================
\section{The Corpus and Ingestion Adapters}
\label{sec:corpus}
% =====================================================================
Table~\ref{tab:corpus} summarises the six datasets and the layer each anchors;
Table~\ref{tab:ingest} reports what was actually ingested. Every adapter emits
schema-valid records, and unmapped source columns are preserved. We do not
re-host raw third-party data; the artifact ships integration scripts and unified
labels, and each source is fetched from its origin.

\begin{table*}[t]
\centering
\caption{The six-dataset corpus. Four network-layer sources and two
navigation-layer sources unify under one schema; only the last row reaches the
mission outcome, and only the C2 source can produce \code{measured\_same\_platform}
pairings.}
\label{tab:corpus}
\small
\renewcommand{\arraystretch}{1.2}
\begin{tabularx}{\textwidth}{@{}l l l l X@{}}
\toprule
\textbf{Dataset} & \textbf{Real/Sim} & \textbf{Layer $\cdot$ sublayer} & \textbf{Event kind} & \textbf{Unique role in the end-to-end benchmark}\\
\midrule
\dataset{UAVIDS-2025}~\cite{uavids2025} & Sim (NS-3.24) & \layernet{} $\cdot$ mesh & flow & Large-scale multi-hop mesh; blackhole/flooding/Sybil/wormhole at scale (122k flows).\\
\dataset{UAV-CAS}~\cite{uavcas2026} & Sim (digital twin) & \layernet{} $\cdot$ mesh & flow & AERPAW-calibrated swarm; collaborative multi-attack compositions (99k flows, 1024 configs).\\
\dataset{UAV Attack Dataset}~\cite{whelan2020uav} & Real testbed & \layernet{} $\cdot$ C2 (MAVLink) & telemetry & Real PX4 GPS spoof/jam + ping-DoS on the \emph{same} airframe family we model---the measured bridge.\\
\dataset{HCRL UAVCAN}~\cite{hcrl_uavcan} & Real testbed & \layernet{} $\cdot$ intra-bus & frame & On-board UAVCAN/DroneCAN flooding/fuzzy/replay---a third attack surface below the C2 link.\\
\dataset{DATAMUt} traces~\cite{datamut2026} & Sim (event) & \layernet{} $\cdot$ mesh (DTN) & hop & Time-window-graph forwarding detector with a \emph{sweepable} operating point $\theta$.\\
\dataset{UAV-EW-Bench-2026}~\cite{uavewbench2026} & Sim (physics-informed) & \layerauto{} $\cdot$ GNSS+ctrl & flight & The navigation anchor: full $0$--$40$\,dB $J/S$ sweep with per-flight MCR outcomes.\\
\bottomrule
\end{tabularx}
\end{table*}

\subsection{Adapter design under source irregularity}
An adapter's job is lossless normalization. Three warrant specifics because
their sources are irregular. The \dataset{UAV Attack Dataset} adapter parses PX4
flight logs, computes horizontal position error against a benign reference (or,
absent one, the flight's own pre-attack median), and emits the
\code{measured\_same\_platform} pairings that ground the interface mapping of
\S\ref{sec:interface}. The \dataset{HCRL UAVCAN} adapter parses candump-style
bus frames; because the per-scenario attack labels are not machine-readable in
the release, we \emph{derive} each scenario's class from the frames
themselves---burst rate relative to the normal bus rate, the fraction of
distinct payloads, and per-frame byte entropy cleanly separate a flooding burst
(one payload hammered above bus rate) from a fuzzy burst (near-random payloads)
from a replay burst (a small set of captured payloads reused).
Table~\ref{tab:hcrl} reports the result, flagged for confirmation against the
dataset's technical report. The \dataset{UAV-CAS} adapter parses the calibrated
digital twin's flow statistics including its stringified list columns; it is
verified against CSVs produced by the original authors' own generator.

\begin{table}[t]
\centering
\caption{Real ingested corpus. Counts are from schema-validated adapter outputs;
every record validates against the schema of \S\ref{sec:schema}.}
\label{tab:ingest}
\small
\renewcommand{\arraystretch}{1.1}
% AUTO-GENERATED by experiments/make_paper_figures.py (clean campaign 2026-07-28).
% PROVENANCE: real-corpus counts from results/ingest_stats.csv (schema-validated staged outputs)
% Do not hand-edit; re-run the script.
\begin{tabularx}{\columnwidth}{@{}Xrr@{}}
\toprule
\textbf{Source (real data ingested)} & \textbf{Events} & \textbf{Windows}\\
\midrule
\dataset{UAV-EW-Bench-2026} & 93,600 & 468\\
\dataset{UAV Attack Dataset} (live sample) & 610 & 0\\
\dataset{UAVIDS-2025} (shard 0) & 122,171 & 4\\
\dataset{HCRL UAVCAN} (10 scenarios) & 200,000 & 34\\
\dataset{UAV-CAS} (stat) & \multicolumn{2}{r}{adapter ready; file pending}\\
\dataset{DATAMUt} traces & \multicolumn{2}{r}{generated per sweep (15{,}392 hops)}\\
\bottomrule
\end{tabularx}

\end{table}

\begin{table}[t]
\centering
\caption{Frame-derived attack classes for the \dataset{HCRL UAVCAN} scenarios,
inferred from burst rate, payload-distinct share, and byte entropy; scenarios
exhibiting more than one burst shape map to \code{collaborative}.}
\label{tab:hcrl}
\small
\renewcommand{\arraystretch}{1.05}
% AUTO-GENERATED by experiments/make_paper_figures.py (clean campaign 2026-07-28).
% PROVENANCE: real-corpus frames; class assignment frame-derived (burst rate, payload-distinct share, byte entropy) pending report confirmation
% Do not hand-edit; re-run the script.
\begin{tabularx}{\columnwidth}{@{}l r r X@{}}
\toprule
\textbf{Scenario} & \textbf{Frames} & \textbf{Bursts} & \textbf{Frame-derived class(es)}\\
\midrule
\code{type1} & 207,858 & 3 & flooding\\
\code{type2} & 134,170 & 3 & replay\\
\code{type3} & 197,479 & 3 & fuzzy\\
\code{type4} & 133,374 & 3 & fuzzy\\
\code{type5} & 180,608 & 3 & fuzzy\\
\code{type6} & 241,321 & 4 & fuzzy\\
\code{type7} & 234,162 & 4 & replay, fuzzy\\
\code{type8} & 265,800 & 4 & fuzzy\\
\code{type9} & 230,378 & 4 & replay, fuzzy\\
\code{type10} & 207,380 & 3 & replay, fuzzy\\
\bottomrule
\end{tabularx}

\end{table}

% =====================================================================
\section{The Detector-Operating-Curve Harness}
\label{sec:harness}
% =====================================================================
A cross-layer guarantee needs the detector's behaviour \emph{as a function of
its operating point}. We expose $\theta$ as a swept parameter and, for each
value, measure detection recall and false-positive rate over multiple seeds,
three network topologies, four routing policies, and two attacker delay regimes.

The resulting operating curve (Fig.~\ref{fig:opcurve}) is the network-side input
a composed guarantee consumes: it tells us, for any target false-positive
budget, how large an undetected per-hop delay an evading attacker may
inject---the residual budget $\Delta(\theta)$ of \eqref{eq:pairing}. The knee in
the curve is structural: an evading attacker's per-hop delay is additive up to
the inter-node contact-window slack, beyond which a missed contact window forces
a full-period penalty.

The knee, and the floor it descends to, are not simulation artifacts but a
structural property of any store-and-forward schedule, and they recur across all
three topologies. Two observations transfer directly to the guarantee. First,
the floor height differs by topology because it equals the probability that a
malicious delay overshoots the slack and self-exposes---higher in the medium
delay regime ($U[1,10]$\,s) than the low one ($U[1,7]$\,s), because a wider
delay draw is likelier to overshoot. Second, no topology's recall falls to zero:
past the knee the detector still catches the self-exposing fraction, which is
precisely why an evading attacker's undetected budget \emph{saturates} rather
than growing without bound. That is the empirical basis for
Lemma~\ref{lem:residual}.

\begin{figure}[t]
\centering
\begin{subfigure}{\columnwidth}
\centering
\begin{tikzpicture}
\begin{axis}[width=\columnwidth, height=3.9cm,
  title={\footnotesize low mode $U[1,7]$\,s},
  xlabel={$\varepsilon$ (s)}, ylabel={recall},
  xmin=0, xmax=10, ymin=0, ymax=1.05,
  legend pos=north east, legend cell align=left, grid=both,
  grid style={black!12}, tick label style={font=\scriptsize},
  label style={font=\scriptsize}, legend style={font=\tiny}]
% AUTO-GENERATED by experiments/make_paper_figures.py (clean campaign 2026-07-28).
% PROVENANCE: simulation-derived (clean campaign, 8 seeds x 4 policies)
% Do not hand-edit; re-run the script.
\addplot[netcol,mark=*,thick] coordinates {(0.25,0.938)(0.5,0.938)(1,0.938)(1.5,0.938)(2,0.938)(3,0.938)(4,0.609)(5,0.234)(6,0.234)(6.5,0.234)(7,0.234)(8,0.234)(10,0.234)};
\addlegendentry{grid $4{\times}4$}
\addplot[bridgecol,mark=square*,thick] coordinates {(0.25,1.000)(0.5,1.000)(1,1.000)(1.5,0.938)(2,0.938)(3,0.938)(4,0.500)(5,0.125)(6,0.125)(6.5,0.125)(7,0.125)(8,0.125)(10,0.125)};
\addlegendentry{ring-9}
\addplot[autocol,mark=diamond*,thick] coordinates {(0.25,1.000)(0.5,1.000)(1,1.000)(1.5,1.000)(2,1.000)(3,0.812)(4,0.562)(5,0.312)(6,0.312)(6.5,0.312)(7,0.312)(8,0.312)(10,0.312)};
\addlegendentry{double ring}

\end{axis}
\end{tikzpicture}
\end{subfigure}

\begin{subfigure}{\columnwidth}
\centering
\begin{tikzpicture}
\begin{axis}[width=\columnwidth, height=3.9cm,
  title={\footnotesize medium mode $U[1,10]$\,s},
  xlabel={$\varepsilon$ (s)}, ylabel={recall},
  xmin=0, xmax=10, ymin=0, ymax=1.05,
  legend pos=north east, legend cell align=left, grid=both,
  grid style={black!12}, tick label style={font=\scriptsize},
  label style={font=\scriptsize}, legend style={font=\tiny}]
% AUTO-GENERATED by experiments/make_paper_figures.py (clean campaign 2026-07-28).
% PROVENANCE: simulation-derived (clean campaign, 8 seeds x 4 policies)
% Do not hand-edit; re-run the script.
\addplot[netcol,mark=*,thick] coordinates {(0.25,0.797)(0.5,0.797)(1,0.797)(1.5,0.797)(2,0.797)(3,0.766)(4,0.750)(5,0.656)(6,0.656)(6.5,0.656)(7,0.656)(8,0.656)(10,0.656)};
\addlegendentry{grid $4{\times}4$}
\addplot[bridgecol,mark=square*,thick] coordinates {(0.25,1.000)(0.5,1.000)(1,1.000)(1.5,1.000)(2,1.000)(3,1.000)(4,0.938)(5,0.625)(6,0.625)(6.5,0.625)(7,0.625)(8,0.625)(10,0.625)};
\addlegendentry{ring-9}
\addplot[autocol,mark=diamond*,thick] coordinates {(0.25,1.000)(0.5,1.000)(1,1.000)(1.5,1.000)(2,1.000)(3,0.875)(4,0.812)(5,0.625)(6,0.625)(6.5,0.625)(7,0.625)(8,0.625)(10,0.625)};
\addlegendentry{double ring}

\end{axis}
\end{tikzpicture}
\end{subfigure}
\caption{The detector operating curve across all three topologies (grid, ring,
double-ring), both delay regimes, $8$ seeds $\times\,4$ routing policies. Recall
holds near its ceiling while $\varepsilon$ stays under the $5$\,s contact-window
slack, then falls onto a floor set by the missed-window self-exposure penalty (a
delay past the slack costs a full TWiG period and is flagged at any
$\varepsilon$). False-positive rate is zero throughout; the measured benign
residual ceiling is $0.178$\,s. The knee-and-floor shape is universal; the floor
height is the topology-specific self-exposure probability.}
\label{fig:opcurve}
\end{figure}

\subsection{Generalising the harness to learned detectors}
\label{sec:roc}
We instantiate the harness on a threshold detector because its operating point
is a single, physically meaningful scalar, which makes the cross-layer sweep
unambiguous. The framework does not depend on that choice, and it is worth
showing precisely why, since most deployed UAV intrusion detection uses learned
classifiers rather than deterministic time-window tests.

A learned detector emits a score $\sigma(a)\in[0,1]$ for a candidate event $a$
and flags it when $\sigma(a)\ge P_{\mathrm{th}}$. Sweeping $P_{\mathrm{th}}$
traces the familiar receiver operating characteristic (ROC): each cutoff is a
$(\mathrm{FPR},\mathrm{TPR})$ pair, exactly as each $\varepsilon$ is in
Fig.~\ref{fig:opcurve}. The only object the composition actually consumes is the
residual budget---the largest attack effect that survives the detector---so the
framework generalises as soon as $P_{\mathrm{th}}$ can be mapped to one. It can,
through the detector's \emph{score-to-effect} profile.

Let $g(\cdot)$ be the monotone-increasing map from an attack's physical
magnitude $\mu$ (here, injected per-hop delay in seconds) to the detector's
expected score, $\sigma=g(\mu)$; $g$ is measured by running the detector over
labelled attacks of known magnitude, which every corpus in \S\ref{sec:corpus}
supports. An attacker evading the cutoff must keep $g(\mu)<P_{\mathrm{th}}$,
hence
\begin{equation}
\label{eq:pth}
\mu \;<\; g^{-1}(P_{\mathrm{th}}) \;=:\; \theta_{\mathrm{eff}}(P_{\mathrm{th}}),
\end{equation}
and Lemma~\ref{lem:residual} applies verbatim with $\theta$ replaced by the
\emph{effective} operating point $\theta_{\mathrm{eff}}$. The threshold detector
is simply the special case $g=\mathrm{id}$, for which
$\theta_{\mathrm{eff}}=\varepsilon$. Two properties of $g$ carry over into the
guarantee. Monotonicity is what makes $g^{-1}$ well defined and the certified
window an interval rather than a union; where a classifier is non-monotone in
$\mu$, one takes the conservative
$\theta_{\mathrm{eff}}=\sup\{\mu: g(\mu)<P_{\mathrm{th}}\}$, which is sound but
looser. And because a learned score is stochastic, $g$ is a distribution rather
than a function, so $\theta_{\mathrm{eff}}$ must be read at an upper confidence
bound on the score---which is precisely the regime in which the
randomized-smoothing certificate of \S\ref{sec:certs}, rather than the
deterministic Gr\"onwall tube, is the appropriate downstream bound.

The practical consequence is that any detector exposing a tunable cutoff and a
measurable score-to-magnitude profile slots into the framework without a schema
change: the harness sweeps $P_{\mathrm{th}}$ instead of $\varepsilon$, records
$(\mathrm{FPR},\mathrm{TPR},\theta_{\mathrm{eff}})$ per cutoff, and the
downstream chain is unchanged. What the framework requires of a detector is not
that it be deterministic, but that its operating point be reducible to a bound
on the physical attack magnitude that survives it.

% =====================================================================
\section{The Interface Mapping $\theta \mapsto \delta(\theta)$}
\label{sec:interface}
% =====================================================================
The interface is where a network-security quantity becomes a navigation-security
quantity. It is the crux of the composition, and---as \S\ref{sec:eval}
shows---the quantity that actually decides whether a useful detector setting is
certifiable.

\subsection{Residual budget of an evading attacker}
\begin{lemma}[Residual delay budget]
\label{lem:residual}
Let $D_\theta$ flag any forwarding hop whose injected delay exceeds $\theta$. An
attacker that evades $D_\theta$ injects at most $\theta$ of undetected delay per
malicious hop---but only up to the contact-window slack $s$: a per-hop delay
exceeding $s$ misses its forwarding window, incurs a full-period penalty
$P_{\!c}$, and is therefore flagged at \emph{any} operating point. Hence
undetected delay is capped at $\min(\theta,s)$ per hop, and over a route with
$H$ attacker-controlled hops the worst-case undetected cumulative delay is
\begin{equation}
\label{eq:delta_budget}
\Delta(\theta) \;\le\; H\,\min(\theta,\,s).
\end{equation}
\end{lemma}
\begin{proof}[Proof sketch]
Evasion requires per-hop delay $\le\theta$. If $\theta\le s$ the delay fits
inside the contact window and the additive term is $H\theta$. If $\theta>s$, any
delay in $(s,\theta]$ overshoots the window, waits a full period
$P_{\!c}\gg s$, and is flagged regardless of $\theta$; an evading attacker
therefore cannot use more than $s$ of invisible delay per hop, giving the $Hs$
cap. Summing over hops yields \eqref{eq:delta_budget}.
\end{proof}

The cap is the \emph{knee}: below the slack the residual budget grows linearly
in $\theta$; above it the budget \emph{saturates} at $Hs$, so the certified
floor plateaus rather than collapsing. We verify the cap directly on the
released harness: over $15{,}392$ observed hops across three topologies, two
attack modes, and eight seeds, no undetected malicious per-hop residual ever
exceeds $4.84$\,s ($s=5$\,s), while every larger delay is flagged as a missed
window ($\ge55$\,s residual). This measurement is what turns the classic
additive budget $H\theta+P_{\!c}\mathbf{1}[\theta>s]$ into the tighter,
saturating \eqref{eq:delta_budget}---an instance of the benchmark improving the
theorem rather than merely testing it.

\subsection{From delay to input perturbation}
A cumulative forwarding delay of $\Delta(\theta)$ seconds means a navigation
correction arrives $\Delta(\theta)$ late; the stack acts on state stale by that
amount. At bounded UAV speed $v_{\max}$, staleness of age $\tau$ bounds the
state/position perturbation:
\begin{equation}
\label{eq:staleness}
\Norm{\delta u} \;\le\; v_{\max}\cdot \Delta(\theta).
\end{equation}
Equation~\eqref{eq:staleness} is the kinematic worst case: it assumes the
aircraft travels at top speed in the wrong direction for the entire staleness
interval. That is sound but pessimistic, because a real estimator does not fly
blind---when GNSS corrections are late or corrupted the EKF down-weights them
and coasts on inertial data, so position error grows far more slowly. We
therefore tighten \eqref{eq:staleness} to
\begin{equation}
\label{eq:gamma}
\Norm{\delta u}\;\le\;\gamma\cdot\Delta(\theta),\qquad \gamma\ll v_{\max},
\end{equation}
where $\gamma$ (metres of position error per second of degradation) is
\emph{measured} on real GPS-spoofing and GPS-jamming flights from the same PX4
airframe family the navigation model represents. Table~\ref{tab:whelan} reports
the calibration: a benign reference flight fixes the pre-attack noise floor, and
the jamming and spoofing flights give the growth rate after onset under two
explicit choices of error signal---the raw receiver fix the attacker injects,
and the EKF-filtered position the controller actually acts on. Because that
choice is itself a certificate-semantics decision we make it an explicit switch
and report both; the certificate consumes the EKF value, which closes the loop.
The gap between the measured $\gamma\approx1.2$--$1.4$\,m/s and
$v_{\max}=15$\,m/s is exactly the slack that decides whether the detector's
operating point is certifiable.

\emph{Temporal validity of the linear model.} Equation~\eqref{eq:gamma} is a
\emph{locally linear} model and must be scoped as one. While GNSS is merely
late rather than absent, the EKF continues to correct against its last accepted
fix and the position error grows approximately linearly in the staleness. Once
the outage exceeds the filter's innovation-gating horizon the estimator is
propagating on inertial data alone, and uncorrected accelerometer bias and
gyroscope drift make the error grow with the square of elapsed time rather than
linearly~\cite{kong2024px4ekf}. We therefore state the model's domain
explicitly: \eqref{eq:gamma} is used only for
$\Delta(\theta)\le\tau_{\mathrm{lin}}$, and beyond $\tau_{\mathrm{lin}}$ a
sound bound requires the second-order form
\begin{equation}
\label{eq:gamma_quad}
\Norm{\delta u}\;\le\;\gamma\,\Delta(\theta)\;+\;\tfrac12\,b_{\mathrm{IMU}}\,\big(\Delta(\theta)-\tau_{\mathrm{lin}}\big)_{+}^{2},
\end{equation}
with $b_{\mathrm{IMU}}$ the residual accelerometer-bias magnitude and
$(\cdot)_+$ the positive part. On the calibration flights the linear fit holds
to within its residual over the full observed degradation interval, giving
$\tau_{\mathrm{lin}}\approx2$\,s; we adopt that value and note that it
exceeds the certified window's upper edge under the measured mapping
($1.33$\,s at a $10$\,m corridor, \S\ref{sec:eval}) by a comfortable margin,
so the entire certified regime lies inside the linear domain and
\eqref{eq:gamma_quad} is never the binding form for the results we report. It
would become binding for a detector tuned near the contact-window slack, where
$\Delta(\theta)$ approaches $Hs=10$\,s---another reason the certified window
is bounded above rather than open.

\begin{table}[t]
\centering
\caption{The measured delay-to-position rate $\gamma$, from three real PX4
flights (one benign reference, one GPS-jamming, one GPS-spoofing) in the
\dataset{UAV Attack Dataset}, self-referenced to the pre-attack median. ``noise''
is the pre-attack position-error $95$th percentile; $\gamma_{\mathrm{lsq}}$ and
$\gamma_{\mathrm{peak}}$ are the least-squares and peak-secant growth rates
after attack onset. These are the \code{measured\_same\_platform} pairings.}
\label{tab:whelan}
\scriptsize
\renewcommand{\arraystretch}{1.1}
% AUTO-GENERATED by experiments/make_paper_figures.py (clean campaign 2026-07-28).
% PROVENANCE: real-corpus (3-flight live sample, self-referenced pre-attack median; hover regime)
% Do not hand-edit; re-run the script.
\begin{tabularx}{\columnwidth}{@{}Xl r r r r@{}}
\toprule
\textbf{Flight} & \textbf{Signal} & \textbf{noise} & \textbf{peak} & \textbf{$\gamma_{\mathrm{lsq}}$} & \textbf{$\gamma_{\mathrm{peak}}$}\\
 & & (m, p95) & (m) & (m/s) & (m/s)\\
\midrule
benign (reference) & receiver & 0.50 & 1.7 & -- & --\\
benign (reference) & ekf & 0.60 & 1.5 & -- & --\\
GPS jamming & receiver & 0.41 & 6.7 & -0.0821 & 0.4805\\
GPS jamming & ekf & 0.47 & 7.0 & -0.1133 & 0.5185\\
GPS spoofing & receiver & 1.14 & 14.3 & 0.2435 & 1.195\\
GPS spoofing & ekf & 0.87 & 15.6 & 0.7007 & 1.365\\
\bottomrule
\end{tabularx}

\end{table}

% =====================================================================
\section{The Composition Theorem}
\label{sec:theorem}
% =====================================================================
\subsection{Navigation dynamics}
Model the UAV navigation state $x(t)\in\Real^n$ as evolving under a (possibly
learned) drift $F$,
\begin{equation}
\label{eq:ode}
\dot{x}(t) = F\big(x(t), u(t)\big), \qquad t\in[0,T],
\end{equation}
where $u(t)$ is the sensor/measurement input over a mission horizon $T$. A
mission \emph{completes} in the conjunctive sense fixed by \eqref{eq:mcr}: the
realised trajectory stays within the safe corridor of margin $m$ for all
$t\in[0,T]$ \emph{and} the objective is reached by the deadline $(1+\kappa)T$.
MCR is the probability of completion over the mission distribution. An attack
perturbs the input, $u\mapsto u+\delta u$, and---for the time-delay class---also
displaces the arrival time.

\begin{theorem}[Composed network-to-navigation guarantee]
\label{thm:main}
Let the navigation drift $F$ of \eqref{eq:ode} be $L$-Lipschitz in its input
over the operating region, with mission horizon $T$ and safe-corridor margin
$m$. Let $D_\theta$ be a network detector with residual budget $\Delta(\theta)$
(Lemma~\ref{lem:residual}) and interface mapping $\delta(\theta)$
(\eqref{eq:gamma} or the kinematic \eqref{eq:staleness}). Then every attack
evading $D_\theta$ induces a trajectory-tube radius bounded by
\begin{equation}
\label{eq:tube}
\rho(\theta) \;=\; \delta(\theta)\,e^{L T},
\end{equation}
and the mission completes---in the conjunctive sense of \eqref{eq:mcr}---whenever
$\rho(\theta)\le m$ and $\Delta(\theta)\le\kappa T$; consequently
\begin{equation}
\label{eq:floor}
\MCR \;\ge\; f\big(\delta(\theta)\big) \;=\;
\Pr_{\text{missions}}\!\big[\,\delta(\theta)\,e^{LT}\le m
\;\wedge\; \Delta(\theta)\le\kappa T\,\big].
\end{equation}
\end{theorem}
\begin{proof}[Proof sketch]
Let $e(t)$ be the deviation between the attacked and nominal trajectories from a
common initial state. $L$-Lipschitzness of $F$ gives
$\Norm{\dot e}\le L\Norm{e}+L\Norm{\delta u}$; Gr\"onwall's
inequality~\cite{gronwall1919} then bounds $\Norm{e(T)}$ by
$\delta(\theta)e^{LT}$, which is \eqref{eq:tube} and discharges the spatial
predicate whenever $\rho(\theta)\le m$. The temporal predicate follows because
the only mechanism by which an evading time-delay attack displaces arrival is
its accumulated undetected delay, bounded by $\Delta(\theta)$ in
Lemma~\ref{lem:residual}. Taking probabilities over the mission distribution
gives \eqref{eq:floor}. Appendix~\ref{app:proof} gives the full argument.
\end{proof}

\emph{On the Lipschitz hypothesis.} Theorem~\ref{thm:main} rests on $F$ being
$L$-Lipschitz \emph{over the operating region}, and that qualifier carries real
physical content rather than being a technical convenience. It asserts that the
aircraft remains in a smooth, locally linearisable flight regime: attitudes away
from gimbal-lock singularities, airspeeds below stall or the onset of blade
stall, actuators off their saturation limits, and no contact dynamics. These
conditions hold for the nominal cruise and loiter envelopes our corpus covers,
and the constant we use is measured inside that envelope rather than assumed
globally (\S\ref{sec:eval}). They can fail: an aggressive evasive manoeuvre, a
stall-recovery transient, or actuator saturation under a strong gust all push
the dynamics into a regime where no finite $L$ bounds $F$ on the relevant set,
and there the certificate is silent rather than merely loose. This is a genuine
scope restriction and we state it as one. It is also the precise point at which
the runtime-assurance reading of \S\ref{sec:discussion} becomes operational: the
certified window is the region in which the learned controller may run, and
departure from the linearisable envelope is exactly the condition a
Simplex-style monitor must treat as a trigger to hand control to the verified
fallback, whose guarantee does not depend on $L$.

\begin{proposition}[Soundness]
\label{prop:sound}
The certified floor never overstates the achievable completion rate: for every
$\theta$, the true MCR of any $D_\theta$-evading attack satisfies
$\MCR\ge f(\delta(\theta))$.
\end{proposition}
\begin{proof}
Every step above is an inequality in the conservative direction: the Lipschitz
bound over-estimates deviation growth, the input bound uses the worst-case
$\delta(\theta)$, and the corridor test is a sufficient---not
necessary---condition for completion. Hence $f(\delta(\theta))$ lower-bounds the
true completion probability. The empirical validation of \S\ref{sec:eval}
confirms the certified floor lies below the measured rate.
\end{proof}

Two consequences frame the rest of the paper. First, \eqref{eq:floor} is
\emph{monotone and then saturating} in $\theta$: tightening the detector shrinks
the residual budget and raises the certified floor, but only until $\theta$
reaches the slack $s$, past which the budget saturates at $Hs$ and the floor
plateaus---a quantitative statement of both how much network security ``buys''
in navigation safety and where that purchase stops. Second, the exponential
$e^{LT}$ in \eqref{eq:tube} is the core certificate's principal limitation: for large
$L$ or long horizons the deterministic bound becomes loose, motivating the
complementary certificates of \S\ref{sec:certs}.

\begin{figure}[t]
\centering
\begin{tikzpicture}[
  font=\normalsize,
  stage/.style={rounded corners=3pt,draw,align=center,minimum height=8.5mm,
    text width=42mm,inner sep=4pt},
  net/.style={stage,fill=netcol!12,draw=netcol,thick},
  br/.style={stage,fill=bridgecol!15,draw=bridgecol,very thick},
  aut/.style={stage,fill=autocol!12,draw=autocol,thick},
  lbl/.style={font=\normalsize\bfseries,fill=white,inner sep=1.5pt},
  >={Latex[length=2.6mm]},
  node distance=8mm]
\node[net] (theta) {detector point $\theta$};
\node[net,below=of theta] (delta0) {residual budget $\Delta(\theta)$};
\node[br,below=of delta0] (delta) {perturbation $\delta(\theta)$};
\node[aut,below=of delta] (tube) {trajectory tube $\delta\,e^{LT}$};
\node[aut,below=of tube] (mcr) {certified floor $\MCR\ge f$};
\draw[->,thick] (theta) -- node[lbl,right] {Lemma~\ref{lem:residual}} (delta0);
\draw[->,thick] (delta0) -- node[lbl,right] {interface (\S\ref{sec:interface})} (delta);
\draw[->,thick] (delta) -- node[lbl,right] {Gr\"onwall} (tube);
\draw[->,thick] (tube) -- node[lbl,right] {corridor margin $m$} (mcr);
\begin{scope}[on background layer]
\node[fit=(theta)(delta0),draw=netcol!55,rounded corners,inner sep=6pt,
  label={[font=\small\bfseries,text=netcol]left:\rotatebox{90}{network side}}]{};
\node[fit=(tube)(mcr),draw=autocol!55,rounded corners,inner sep=6pt,
  label={[font=\small\bfseries,text=autocol]left:\rotatebox{90}{autonomy side}}]{};
\end{scope}
\end{tikzpicture}
\caption{The composition chain of Theorem~\ref{thm:main}. A detector operating
point flows top-to-bottom through a residual-budget bound, an empirically
grounded interface mapping, a Gr\"onwall trajectory tube, and a corridor-margin
test, arriving at a certified mission-completion floor. The upper two stages are
network-side quantities, the lower two autonomy-side; $\delta(\theta)$ is the
bridge between them.}
\label{fig:chain}
\end{figure}

% =====================================================================
\section{A Four-Certificate Architecture}
\label{sec:certs}
% =====================================================================
Theorem~\ref{thm:main} uses the Lipschitz--Gr\"onwall certificate as its spine.
Does that single certificate suffice? It does not. The bound is
\emph{deterministic} and \emph{worst-case}: it holds for any perturbation up to
its radius, but the radius collapses as $LT$ grows, and it presumes a Lipschitz
constant that is tractable to bound. Three regimes escape it, each covered by a
distinct certificate (Table~\ref{tab:certs}).

\begin{table}[t]
\centering
\caption{The four certificates cover disjoint regimes; none subsumes another.
The Lipschitz--Gr\"onwall core gives the deterministic MCR floor; the other
three extend coverage where it is loose or silent.}
\label{tab:certs}
\small
\renewcommand{\arraystretch}{1.15}
\begin{tabularx}{\columnwidth}{@{}
  >{\hsize=0.82\hsize\raggedright\arraybackslash}X
  >{\hsize=0.70\hsize\raggedright\arraybackslash}X
  >{\hsize=1.48\hsize\raggedright\arraybackslash}X@{}}
\toprule
\textbf{Certificate} & \textbf{Guarantee type} & \textbf{Unique niche}\\
\midrule
Lipschitz--Gr\"onwall & deterministic, worst-case & The MCR floor: bounded input $\Rightarrow$ bounded tube $\Rightarrow$ bounded loss. Tight only for moderate $LT$.\\
Randomized smoothing~\cite{cohen2019smoothing} & probabilistic radius & Scales to large / non-Lipschitz-tractable models; certifies when the budget is known only statistically.\\
PAC-Bayes~\cite{viallard2021pacbayes} & distributional & Transfers the floor from the training mission distribution to unseen deployment regions.\\
MWU regret~\cite{arora2012mwu,balcan2015commitment} & online, adaptive & Bounds loss against an \emph{adaptive} jammer over time---a dimension no fixed-ball certificate addresses.\\
\bottomrule
\end{tabularx}
\end{table}

\paragraph{Why the core does not subsume the others}
The Lipschitz--Gr\"onwall radius is a single deterministic ball; randomized
smoothing certifies a probabilistic ball for models whose Lipschitz constant is
intractable, trading determinism for scalability---a strictly different
operating point, not a special case. PAC-Bayes certifies a \emph{distribution}
of models over a \emph{distribution} of missions, orthogonal to any per-input
radius. MWU regret certifies behaviour against an adversary that \emph{adapts
over time}, which no static perturbation ball can express. The four form a
cover, not a hierarchy.

For concreteness we record the two complementary bounds we instantiate
numerically. Randomized smoothing~\cite{cohen2019smoothing} wraps the base model
in Gaussian noise of level $\sigma$ and certifies, with confidence $1-\alpha$
over $n$ samples, an $\ell_2$ radius
\begin{equation}
\label{eq:cohen}
R_{\mathrm{RS}} \;=\; \tfrac{\sigma}{2}\big(\Phi^{-1}(\underline{p_A})-\Phi^{-1}(\overline{p_B})\big),
\end{equation}
with $\Phi^{-1}$ the inverse standard-normal CDF and
$\underline{p_A},\overline{p_B}$ confidence bounds on the top-two class
probabilities. The PAC-Bayes
complement~\cite{mcallester1999pac,viallard2021pacbayes} transfers a bound from
the training mission distribution to deployment: with empirical risk
$R_{\mathrm{emp}}$ over $m$ missions and prior-to-posterior divergence
$\mathrm{KL}$, with probability $1-\delta$,
\begin{equation}
\label{eq:mcallester}
R_{\mathrm{true}} \;\le\; R_{\mathrm{emp}} + \sqrt{\tfrac{\mathrm{KL}+\ln(2\sqrt{m}/\delta)}{2m}}.
\end{equation}
Our engine evaluates \eqref{eq:mcallester} against the trained model's measured
$R_{\mathrm{emp}}$; the single quantity it awaits is the numeric $\mathrm{KL}$,
which we flag as pending rather than substitute (Table~\ref{tab:constants}).

\emph{How the four operate on a single flight.} The four certificates are
concurrent, not alternative, and the clearest way to see that is to hold one
flight fixed. Consider a single ten-minute delivery sortie, flown by a model
trained on last season's mission mix, through a corridor in which a compromised
relay delays peer corrections and a ground jammer varies its transmit power.
Four claims are in force simultaneously, and they are claims about four
different objects. (i) The Lipschitz--Gr\"onwall tube bounds the trajectory
deviation produced by the \emph{deterministic, bounded} component of the
disturbance---the undetected forwarding delay $\Delta(\theta)$---and it is the
only one of the four that yields a per-instant geometric envelope, which is why
it, and it alone, discharges the corridor test of \eqref{eq:mcr}. (ii)
Randomized smoothing bounds the effect of the \emph{stochastic} component: the
jammer's residual perturbation of the receiver feature vector admits no hard
bound, only a distribution, and the resulting claim is about the model's
\emph{decision} at a point rather than about the trajectory. (iii) PAC-Bayes
bounds the gap between the risk this model exhibited on its training mission mix
and the risk it exhibits on \emph{this} sortie, whose parameters were not in
that mix; it is a statement about the model, constant across the flight rather
than attached to any instant of it. (iv) MWU regret bounds the cumulative loss
of the defender's \emph{policy switching} over the sortie against a jammer that
adapts; it is a statement about the time axis and is vacuous at any single
instant.

Because a trajectory, a decision, a model, and a policy sequence are distinct
objects, the four compose by conjunction rather than selection: a deployment
reads its floor as the pointwise minimum
$\min\{f_{\mathrm{G}},f_{\mathrm{RS}},f_{\mathrm{PB}},f_{\mathrm{MWU}}\}$ at its
operating point, and the binding term identifies which assumption is closest to
failing. In the regime this paper evaluates---a bounded time-delay disturbance,
in-distribution missions, a non-adaptive adversary---the Gr\"onwall term binds
and the other three are slack, which is why the headline result of
\S\ref{sec:eval} is a Gr\"onwall curve. We nonetheless carry all four, because
relaxing any one of those three conditions moves the binding term, and a
framework that went silent at that point would not be an end-to-end guarantee.
We are correspondingly explicit about maturity: the Gr\"onwall and
randomized-smoothing floors are instantiated numerically on the reference model,
the MWU regret bound is instantiated at $|S|=4$, and the PAC-Bayes bound is
reported as a form awaiting its numeric $\mathrm{KL}$ rather than as a number
(Table~\ref{tab:constants}).

% =====================================================================
\section{Instantiation and Evaluation}
\label{sec:eval}
% =====================================================================
\subsection{The navigation anchor}
The autonomy layer supplies the outcome the whole construction reaches for.
Fig.~\ref{fig:mcr} shows MCR versus $J/S$ for four defence configurations,
computed from the released benchmark's $93{,}600$-flight per-flight table. The
undefended baseline collapses below the $0.90$ floor by $J/S\!\approx\!8$\,dB,
while the strongest configuration holds it to $J/S\!\approx\!25$\,dB. For our
purpose the point is not which defence wins but that the autonomy layer produces
a \emph{graded, measurable} outcome a network event can be paired against.

\begin{figure}[t]
\centering
\begin{tikzpicture}
\begin{axis}[
  width=\columnwidth, height=4.2cm,
  xlabel={jamming-to-signal ratio $J/S$ (dB)},
  ylabel={Mission Completion Rate},
  xmin=0, xmax=40, ymin=0, ymax=1.03,
  legend pos=south west, legend cell align=left,
  grid=both, grid style={black!12},
  tick label style={font=\footnotesize},
  label style={font=\footnotesize},
  legend style={font=\scriptsize}]
% AUTO-GENERATED by experiments/make_paper_figures.py (clean campaign 2026-07-28).
% PROVENANCE: real-corpus (93,600-flight per_flight.csv, Wilson CIs in results/ewbench_mcr_anchor.csv)
% Do not hand-edit; re-run the script.
\addplot[accent,mark=triangle*,thick] coordinates {(0,0.990)(5,0.947)(10,0.832)(15,0.510)(20,0.263)(25,0.125)(30,0.038)(35,0.013)(40,0.000)};
\addlegendentry{No-Def (PX4)}
\addplot[bridgecol,mark=square*,thick] coordinates {(0,0.987)(5,0.955)(10,0.917)(15,0.847)(20,0.692)(25,0.510)(30,0.323)(35,0.178)(40,0.068)};
\addlegendentry{CAF-CNN}
\addplot[autocol,mark=diamond*,thick] coordinates {(0,1.000)(5,0.970)(10,0.948)(15,0.913)(20,0.763)(25,0.635)(30,0.510)(35,0.275)(40,0.120)};
\addlegendentry{Seq2Seq Tr.}
\addplot[netcol,mark=*,thick] coordinates {(0,1.000)(5,1.000)(10,0.998)(15,0.978)(20,0.938)(25,0.913)(30,0.793)(35,0.710)(40,0.530)};
\addlegendentry{Ours (M1+M4+M6+M7)}

\draw[black,dashed] (axis cs:0,0.90) -- (axis cs:40,0.90);
\node[anchor=west,font=\scriptsize] at (axis cs:0.5,0.855) {DO-326A floor 0.90};
\end{axis}
\end{tikzpicture}
\caption{The navigation anchor: MCR vs.\ $J/S$ on \dataset{UAV-EW-Bench-2026},
computed from the released $93{,}600$-flight table (Wilson $95\%$ intervals in
the artifact). The graded outcome is what a network attack window is paired
against.}
\label{fig:mcr}
\end{figure}

\subsection{Reference constants}
We instantiate the certificate on the reference navigation model, a
continuous-time graph model of the GNSS constellation and receiver. Global
power-iteration gives a Lipschitz estimate $\hat L\approx1.01$; measuring the
constant \emph{along the hidden-state trajectories the integrator actually
visits on operating-region inputs} gives a larger local value
$\hat L_{\mathrm{loc}}\approx1.18$. We report the tighter local radius as the
defensible number: with $T=1$ and $\varepsilon_{\mathrm{out}}=0.5$,
$R=\varepsilon_{\mathrm{out}}e^{-\hat L_{\mathrm{loc}}T}\approx0.153$ (the
global estimate gives $0.182$). Note the local radius is \emph{smaller}: the
bound is not loose in the direction that would flatter us, and the certified
regime is governed by real sensitivity rather than an optimistic estimate. The
smoothing complement gives a certified $\ell_2$ radius $\approx0.44$ at
$\sigma=0.25$ ($\alpha=10^{-3}$, $n=200$; measured mean $0.446$ on the trained
checkpoint), and the MWU defender mixes four policies with regret
$R(T)\le\sqrt{T\ln|S|}$, $|S|=4$. Constants come from the trained reference
implementation on a transparent synthetic corpus and are held fixed;
Table~\ref{tab:constants} lists each with its provenance tag.

\begin{table}[t]
\centering
\caption{Instantiated constants and their provenance. The one pending numeric
constant is the PAC--Bayes prior/posterior KL, which we flag rather than
substitute.}
\label{tab:constants}
\small
\renewcommand{\arraystretch}{1.15}
% AUTO-GENERATED by experiments/make_paper_figures.py (clean campaign 2026-07-28).
% PROVENANCE: mixed; every row carries its own tag
% Do not hand-edit; re-run the script.
\begin{tabularx}{\columnwidth}{@{}p{1.9cm}@{\hspace{5pt}}X@{\hspace{5pt}}p{1.75cm}@{}}
\toprule
\textbf{Component} & \textbf{Constant} & \textbf{Provenance}\\
\midrule
Lipschitz--Gr\"onwall & $\hat L_{\mathrm{loc}}{=}1.181\Rightarrow R{=}0.153$; global $\hat L{=}1.013\Rightarrow R{=}0.182$ & trained ckpt [fixture]\\
Rand.\ smoothing & $\sigma{=}0.25$, $R_{\ell_2}{=}0.44$ (measured $0.446$) & trained ckpt [fixture]\\
PAC--Bayes & McAllester on real $R_{\mathrm{emp}}$; numeric KL pending & methods ch.\ [pending]\\
MWU regret & $R(T)\le\sqrt{T\ln 4}$; $T{=}100\Rightarrow11.77$ & exact\\
Unit bridge (RF classes) & \code{caf\_shift\_v2}: $0.45$\,m (Gr\"onwall) / $1.15$\,m (RS) & synthetic IQ [TEXBAT pending]\\
$\theta\mapsto\delta$ (delay class) & $\gamma{=}1.20$--$1.37$\,m/s measured; kinematic $15$\,m/s fallback & real 3 flights [hover]\\
\bottomrule
\end{tabularx}

\end{table}

\subsection{A worked end-to-end example}
\label{sec:worked}
To make \eqref{eq:chain} concrete we trace one number through it. Take the grid
topology, whose routes carry $H=2$ attacker-controlled hops, and the detector
setting $\theta=0.25$\,s. \emph{(1) Threshold to budget.} Since
$\theta<s=5$\,s, Lemma~\ref{lem:residual} gives $\Delta(\theta)\le H\theta=0.50$\,s
of undetected staleness. \emph{(2) Budget to perturbation.} With the measured
EKF rate $\gamma=1.37$\,m/s, \eqref{eq:gamma} gives
$\Norm{\delta u}\le0.69$\,m; the kinematic bound would instead give $7.5$\,m.
\emph{(3) Perturbation to tube.} With $\hat L_{\mathrm{loc}}=1.18$ and $T=1$,
the tube radius is $\rho=0.69\,e^{1.18}\approx2.25$\,m. \emph{(4) Tube to
floor.} For any corridor margin $m\ge2.25$\,m the mission is certified to
complete; across the margin family $\{2,5,10,20\}$\,m the certified floor is
positive for all but the tightest $2$\,m corridor. Under the kinematic bound,
step (3) yields $\rho\approx24.5$\,m, exceeding every margin---which is exactly
why the empirical calibration, not the theorem, decides the framing.
Fig.~\ref{fig:budget} shows the measured undetected budget feeding step (2),
confirming the saturation predicted by Lemma~\ref{lem:residual}. Step (5), the
temporal predicate of \eqref{eq:mcr}, is satisfied by a wide margin at this
operating point: the induced schedule slip is $\Delta(\theta)\le0.50$\,s, which
is under a $10\%$ tolerance for any mission longer than five seconds.

\emph{Why these corridor margins.} The family $\{2,5,10,20\}$\,m is not
arbitrary; it brackets the lateral buffers that operational approval frameworks
already require, which is what makes the certified window comparable to a
regulatory quantity rather than to a modelling choice. Under the JARUS
Specific Operations Risk Assessment (SORA), the reference methodology adopted by
EASA and other authorities for approving UAS operations in the
\emph{specific} category, the operational volume comprises the flight geography
plus a \emph{contingency volume} whose minimum lateral extent is the sum of
independent error terms: GNSS position accuracy, position-holding error, map
error, a reaction-distance term, and the stopping distance of the contingency
manoeuvre~\cite{jarus_sora25,easa_uspace2021}. SORA v2.5's default values for
the first three alone---$3$\,m GNSS accuracy, $3$\,m position-holding error,
$1$\,m map error---already total $7$\,m before any reaction or manoeuvre term is
added. Our margin family therefore spans the regulator's own scale: $2$\,m is
\emph{tighter} than any SORA-compliant lateral buffer and functions as a
deliberately adversarial stress case; $5$\,m sits inside the default static
budget; $10$\,m is representative of a small multirotor's full contingency
extent once reaction distance is included; and $20$\,m covers a fixed-wing
platform, whose larger turn radius and higher cruise speed inflate the
manoeuvre term. Reading Fig.~\ref{fig:sensitivity} against these values converts
the abstract question ``is $\theta$ certifiable?'' into the operational one an
approval submission must answer: whether the detector setting keeps the
worst-case evading tube inside the contingency volume the operator has already
declared.

\begin{figure}[t]
\centering
\begin{tikzpicture}
\begin{axis}[
  width=\columnwidth, height=4.2cm,
  xlabel={detector threshold $\theta$ (s)},
  ylabel={mean malicious delay (s)},
  xmin=0, xmax=10, ymin=0, ymax=26,
  legend cell align=left, legend columns=2,
  legend style={at={(0.5,1.06)}, anchor=south, font=\scriptsize, draw=black!30},
  grid=both, grid style={black!12},
  tick label style={font=\footnotesize},
  label style={font=\footnotesize}]
% AUTO-GENERATED by experiments/make_paper_figures.py (clean campaign 2026-07-28).
% PROVENANCE: simulation-derived (hop ledger means over 8 seeds x 4 policies, scenario 1)
% Do not hand-edit; re-run the script.
\addplot[netcol,mark=*,thick] coordinates {(0.25,0.00)(0.5,0.00)(1,0.00)(1.5,0.00)(2,0.00)(3,0.00)(4,2.40)(5,5.51)(6,5.51)(6.5,5.51)(7,5.51)(8,5.51)(10,5.51)};
\addlegendentry{undetected, low mode}
\addplot[netcol,dashed] coordinates {(0.25,31.60)(0.5,31.60)(1,31.60)(1.5,31.60)(2,31.60)(3,31.60)(4,31.60)(5,31.60)(6,31.60)(6.5,31.60)(7,31.60)(8,31.60)(10,31.60)};
\addlegendentry{injected total, low mode}
\addplot[bridgecol,mark=square*,thick] coordinates {(0.25,0.00)(0.5,0.00)(1,0.00)(1.5,0.00)(2,0.00)(3,0.16)(4,0.28)(5,1.13)(6,1.13)(6.5,1.13)(7,1.13)(8,1.13)(10,1.13)};
\addlegendentry{undetected, medium mode}
\addplot[bridgecol,dashed] coordinates {(0.25,195.89)(0.5,195.89)(1,195.89)(1.5,195.89)(2,195.89)(3,195.89)(4,195.89)(5,195.89)(6,195.89)(6.5,195.89)(7,195.89)(8,195.89)(10,195.89)};
\addlegendentry{injected total, medium mode}
\draw[black,dotted,thick] (axis cs:0.25,0) -- (axis cs:0.25,25);
\draw[black!60,dashed] (axis cs:0.05,10) -- (axis cs:10,10) node[pos=0.02,anchor=south west,font=\small] {$H\cdot s = 10$\,s saturation cap};

\end{axis}
\end{tikzpicture}
\caption{The measured residual budget behind Lemma~\ref{lem:residual}
(scenario~1, mean over $8$ seeds $\times\,4$ policies). The \emph{undetected}
delay an evader slips past the detector (solid) saturates well below the
\emph{total} injected delay (dashed): beyond the contact-window slack the extra
delay produces missed windows flagged at any $\theta$, so the curve the
certificate consumes flattens rather than growing.}
\label{fig:budget}
\end{figure}

\subsection{The certified floor and the interface that decides it}
Fig.~\ref{fig:headline} is the headline result, and it is deliberately
\emph{parametric in the interface mapping}. Under the kinematic worst case,
$\theta=0.25$\,s falls just outside the certified window; under the empirically
measured mapping it falls \emph{inside} at every corridor margin.
Fig.~\ref{fig:sensitivity} quantifies how much room there is: the mapping would
have to steepen past $\gamma_{\mathrm{req}}=7.28$\,m/s (at a $10$\,m corridor)
to push the operating point back out---more than five times the measured rate.
This is the paper's most consequential finding for the field: what determines
whether cross-layer certification is useful is a \emph{measurement} of how
attacks propagate into state error, not the proof machinery.

\begin{figure}[t]
\centering
\begin{tikzpicture}
\begin{semilogxaxis}[
  width=\columnwidth, height=4.9cm,
  xlabel={detector operating point $\theta$ (s, log scale)},
  ylabel={certified floor},
  xmin=0.05, xmax=10, ymin=0, ymax=1.05,
  legend cell align=left, legend columns=2,
  legend style={at={(0.5,1.06)}, anchor=south, font=\scriptsize, draw=black!30,
    inner sep=2pt, nodes={inner xsep=2pt},
    /tikz/every even column/.append style={column sep=0.4em}},
  grid=both, grid style={black!12},
  tick label style={font=\footnotesize},
  label style={font=\footnotesize}]
% AUTO-GENERATED by experiments/make_paper_figures.py (clean campaign 2026-07-28).
% PROVENANCE: real-corpus delta mapping + simulation-derived budgets + stated margin family
% Do not hand-edit; re-run the script.
\fill[netcol!9] (axis cs:0.178,0) rectangle (axis cs:1.3341,1.05);
\addplot[netcol,thick] coordinates {(0.05,1.000)(0.1,1.000)(0.121,1.000)(0.15,1.000)(0.178,1.000)(0.2,1.000)(0.243,1.000)(0.25,1.000)(0.3,0.750)(0.4,0.750)(0.5,0.750)(0.67,0.500)(1,0.500)(1.5,0.250)(2,0.250)(3,0.000)(4,0.000)(5,0.000)(6,0.000)(6.5,0.000)(7,0.000)(8,0.000)(10,0.000)};
\addlegendentry{empirical $\gamma{=}1.37$\,m/s (Whelan EKF)}
\addplot[autocol,thick,dotted] coordinates {(0.05,1.000)(0.1,1.000)(0.121,1.000)(0.15,1.000)(0.178,1.000)(0.2,1.000)(0.243,1.000)(0.25,1.000)(0.3,1.000)(0.4,0.750)(0.5,0.750)(0.67,0.750)(1,0.500)(1.5,0.500)(2,0.250)(3,0.250)(4,0.000)(5,0.000)(6,0.000)(6.5,0.000)(7,0.000)(8,0.000)(10,0.000)};
\addlegendentry{empirical $\gamma{=}1.20$\,m/s (Whelan receiver)}
\addplot[accent,dashed] coordinates {(0.05,0.750)(0.1,0.500)(0.121,0.500)(0.15,0.250)(0.178,0.250)(0.2,0.250)(0.243,0.000)(0.25,0.000)(0.3,0.000)(0.4,0.000)(0.5,0.000)(0.67,0.000)(1,0.000)(1.5,0.000)(2,0.000)(3,0.000)(4,0.000)(5,0.000)(6,0.000)(6.5,0.000)(7,0.000)(8,0.000)(10,0.000)};
\addlegendentry{kinematic worst case $\gamma{=}15$\,m/s}
\addplot[black,thick] coordinates {(0.05,0)(10,0)};
\addlegendentry{autonomy-/network-only floor ($=0$)}
\draw[black,dotted] (axis cs:0.25,0) -- (axis cs:0.25,1.05);
\node[anchor=south,font=\scriptsize,rotate=90] at (axis cs:0.25,0.30) {$\theta{=}0.25$\,s};

\end{semilogxaxis}
\end{tikzpicture}
\caption{The headline result: certified floor vs.\ detector operating point
$\theta$ for the time-delay class (grid topology, Gr\"onwall tube at the measured
local $L{=}1.181$, $H{=}2$, corridor-margin family $\{2,5,10,20\}$\,m). Under
the kinematic worst case (red) $\theta{=}0.25$\,s is just outside the certified
window; under the empirically measured mapping from real GPS-spoof/jam flights
($\gamma{\approx}1.2$--$1.4$\,m/s) it is \emph{inside} at every margin (shaded:
the EKF window $[0.178,1.334]$\,s at $m{=}10$\,m). Both single-layer baselines
certify a zero floor for an evading attacker (black).}
\label{fig:headline}
\end{figure}

\begin{figure}[t]
\centering
\begin{tikzpicture}
\begin{loglogaxis}[
  width=\columnwidth, height=4.2cm,
  xlabel={corridor margin $m$ (m)},
  ylabel={window edge $\theta^{*}(m)$ (s)},
  xmin=1.8, xmax=22, ymin=0.015, ymax=4,
  xtick={2,5,10,20}, xticklabels={2,5,10,20},
  legend cell align=left, legend columns=1,
  legend style={at={(0.5,1.06)}, anchor=south, font=\scriptsize, draw=black!30},
  grid=both, grid style={black!12},
  tick label style={font=\footnotesize},
  label style={font=\footnotesize}]
% AUTO-GENERATED by experiments/make_paper_figures.py (clean campaign 2026-07-28).
% PROVENANCE: real-corpus gamma (Whelan 3-flight hover) + simulation-derived H, slack; theta* = m/(gamma e^{LT} H)
% Do not hand-edit; re-run the script.
\addplot[autocol,mark=diamond*,thick,dotted] coordinates {(2,0.2569)(5,0.6422)(10,1.2844)(20,2.5688)};
\addlegendentry{empirical $\gamma{=}1.20$\,m/s (receiver)}
\addplot[netcol,mark=*,thick] coordinates {(2,0.2249)(5,0.5622)(10,1.1244)(20,2.2489)};
\addlegendentry{empirical $\gamma{=}1.37$\,m/s (EKF)}
\addplot[accent,mark=triangle*,thick,dashed] coordinates {(2,0.0205)(5,0.0512)(10,0.1023)(20,0.2046)};
\addlegendentry{kinematic $\gamma{=}15$\,m/s}
\draw[black,dotted,thick] (axis cs:2,0.25) -- (axis cs:20,0.25) node[pos=0.04,anchor=south west,font=\small] {paper point $\theta{=}0.25$\,s};
\draw[black!60,dashed] (axis cs:2,0.178) -- (axis cs:20,0.178) node[pos=0.04,anchor=north west,font=\small] {benign ceiling $0.178$\,s};

\end{loglogaxis}
\end{tikzpicture}
\caption{The decisive sensitivity: the loosest detector setting
$\theta^{*}(m)=m/(\gamma e^{LT}H)$ whose worst-case evading tube still fits a
corridor of margin $m$, per interface mapping. A mapping certifies the paper
operating point iff its curve lies above the $\theta{=}0.25$\,s line; the
feasible region is bounded below by the measured benign residual ceiling
($0.178$\,s, the FPR$=$0 floor). The measured rates clear the requirement
$\gamma_{\mathrm{req}}{=}7.28$\,m/s at $m{=}10$\,m by ${\sim}5\times$; the
kinematic worst case fails at every margin.}
\label{fig:sensitivity}
\end{figure}

\subsection{Dominance over single-layer baselines}
We state the dominance claim precisely rather than as blanket range-coverage.
The composed guarantee dominates two baselines: (i) \emph{autonomy-only}---a
navigation certificate with no detector, which must assume the full unmitigated
attack including the missed-window penalties (measured cumulative delay up to
$424$\,s) and therefore certifies a zero floor against an evading attacker; and
(ii) \emph{network-only}---a detector with no navigation guarantee, for which
any evading attack leaves the aircraft undefended, floor zero by definition.
\textbf{The composed guarantee is the only configuration with a nonzero
certified floor at all.} Quantifying per seed with the Wilcoxon signed-rank test
and Holm correction, the composed floor exceeds the autonomy-only floor at all
$13$ operating points on the $\theta$ grid, every adjusted $p<10^{-20}$. The
composition wins because it alone uses the detector to \emph{shrink} the
perturbation the certificate must absorb, and---via the tightened budget of
Lemma~\ref{lem:residual}---the floor saturates rather than collapsing past the
knee.

The honest scope: the certified operating window is nonempty but bounded. Under
the empirically calibrated interface it is $[0.178,\,1.33]$\,s at a $10$\,m
corridor, with the operating point $\theta=0.25$\,s comfortably inside; under
the conservative kinematic mapping it shrinks to $[0.178,\,0.243]$\,s and
$\theta=0.25$\,s falls just outside. We claim the former and show the latter as
the worst-case fallback; we do not claim coverage across the entire relaxation
range.

\subsection{Certified-versus-empirical validation, and the two floors}
The certificate is sound iff the certified floor never exceeds the realised MCR
(Proposition~\ref{prop:sound}). At the anchor $J/S=20$\,dB the strongest defence
attains empirical $\MCR=0.938$ over $93{,}600$ real benchmark flights,
comfortably above the certified target $0.80$; the gap is the conservativeness
of the worst-case Gr\"onwall bound, and it narrows where the empirically
tightened interface replaces the kinematic worst case.

We keep two floors distinct throughout, because conflating them would overstate
the guarantee. The reference certificate targets a \emph{certified}
$\MCR\ge0.80$ at the anchor operating point; the empirical benchmark reports the
DO-326A $0.90$ \emph{operational} floor holding further into the band. These are
different quantities---a conservative proof-derived bound versus a measured
rate---and we label every reported number accordingly.

\begin{table}[t]
\centering
\caption{A representative cross-layer pairing: a detector-evading time-delay
attack still induces measurable navigation staleness. Values from the released
artifact (harness simulation; three-flight $\delta$ calibration; benchmark MCR).}
\label{tab:pairing}
\small
\renewcommand{\arraystretch}{1.15}
% AUTO-GENERATED by experiments/make_paper_figures.py (clean campaign 2026-07-28).
% PROVENANCE: simulation-derived (DATAMUt) + real-corpus (Whelan gamma, EW-Bench MCR)
% Do not hand-edit; re-run the script.
\begin{tabularx}{\columnwidth}{@{}lX@{}}
\toprule
\textbf{Quantity} & \textbf{Value at representative $\theta$}\\
\midrule
Detector threshold $\theta$ & $0.25$\,s (paper operating point)\\
Detector verdict on tuned attacker & evaded (by construction)\\
Malicious hops $H$ & $2$ (grid/ring/double-ring routes)\\
Residual budget $\Delta(\theta)\le H\min(\theta,s)$ & $\le 0.50$\,s ($s{=}5$\,s; measured cap $4.84$\,s)\\
Measured undetected delay (192 runs) & median $0.00$\,s, max $0.00$\,s\\
Staleness, empirical $\gamma{=}1.37$\,m/s & $\le 0.69$\,m (Whelan, hover)\\
Staleness, kinematic $v_{\max}{=}15$\,m/s & $\le 7.5$\,m (worst case)\\
Paired autonomy window & \dataset{UAV-EW-Bench} at $J/S{=}20$\,dB\\
Empirical MCR of paired window & $0.938$ (strongest defence)\\
\code{pairing\_basis} & synthetic / measured (per record)\\
\bottomrule
\end{tabularx}

\end{table}

\subsection{Three evidence classes, one worked pairing each}
Table~\ref{tab:pairing} reports the pairing at a representative operating point.
The \code{pairing\_basis} flag is easiest to understand through one example of
each class. \emph{Measured, same platform:} the \dataset{UAV Attack Dataset}
records a real GPS-jamming flight and its measured position error on one
airframe; a pairing joining that flight to the corresponding RF event is tagged
\code{measured\_same\_platform}, and its $\delta$ is a measurement, not a model.
\emph{Physics model:} a forwarding-attack window carries a residual budget in
seconds; coupling it to a navigation window through the stated relation
\eqref{eq:gamma} yields a \code{physics\_model} pairing---a named, versioned
physical relation a reviewer can inspect and challenge, and the class the
certificate consumes. \emph{Synthetic alignment:} when a network window from one
simulated source and an autonomy window from another are aligned by
construction, the pairing is tagged \code{synthetic\_alignment}: maximal
coverage, weakest causal claim, and the flag says so. A user needing a strong
claim filters to the first class; a user needing breadth accepts the third with
its provenance explicit.

\subsection{How much of the benchmark rests on measurement}
\label{sec:provenance_dist}
The flag is only useful if its distribution is published, so we report it rather
than describe it. Table~\ref{tab:provenance} gives the acquisition split across
the whole corpus: $200{,}610$ of $431{,}773$ ingested events ($46.5\%$) were
captured on real hardware, the remainder generated by simulators of three
different kinds. That split is favourable, but it is also the \emph{weaker} of
the two questions a reviewer should ask, and conflating them would overstate the
benchmark.

The stronger question is how many \emph{pairings} rest on measurement, and here
the honest answer is more restrictive. A \code{measured\_same\_platform} pairing
requires one source to have observed both the network event and the flight it
degraded on a single airframe. Exactly one source in the corpus can do this---the
\dataset{UAV Attack Dataset}---and although its three flights supply the
delay-to-position rate $\gamma$ that the certificate consumes
(Table~\ref{tab:whelan}, a real-corpus measurement obtained by
self-referencing each flight to its own pre-attack median), the release does not
carry machine-readable attack on/off timestamps. Without them the adapter cannot
emit the \code{AttackWindow} boundaries a pairing record requires, and we decline
to infer them. \textbf{The number of released \code{measured\_same\_platform}
pairings is therefore currently zero}: every released pairing is
\code{physics\_model} (the network window coupled to a navigation window through
the stated relation \eqref{eq:gamma}, and the class the theorem consumes) or
\code{synthetic\_alignment} (aligned by construction under
Appendix~\ref{app:alignment}, with the error budget \eqref{eq:align_err}).

We state this plainly because it is exactly the fact the flag exists to expose,
and because it locates the benchmark's binding limitation precisely: not the
volume of data, of which there is a great deal, but the scarcity of captures
that observe both layers at once. The headline guarantee does not depend on
closing that gap---$\gamma$ is measured on real flights regardless of whether the
window boundaries are machine-readable---but the strongest evidence class the
schema can express is, today, unpopulated, and a reader filtering the artifact to
\code{measured\_same\_platform} will correctly retrieve nothing. Obtaining
same-platform captures with published attack intervals, across flight regimes
rather than hover alone, is consequently the highest-value extension of the
corpus, and \S\ref{sec:limitations} treats it as such.

\begin{table}[t]
\centering
\caption{Acquisition provenance across the ingested corpus. Counts are from the
schema-validated adapter outputs; the split is machine-generated by
\code{experiments/provenance\_ledger.py} so the paper cannot drift from the
artifact. Acquisition is a property of the \emph{source}; the strictly stronger
\code{pairing\_basis} distribution is discussed in the text.}
\label{tab:provenance}
\small
\renewcommand{\arraystretch}{1.15}
% MACHINE-GENERATED by experiments/provenance_ledger.py -- do not edit.
% Regenerate after any adapter run; the paper must not drift from the artifact.
\begin{tabularx}{\columnwidth}{@{}l l r r@{}}
\toprule
\textbf{Source} & \textbf{Acquisition} & \textbf{Events} & \textbf{\%}\\
\midrule
\dataset{hcrl\_uavcan} & real testbed & 200,000 & 46.3\\
\dataset{uavids\_2025} & simulation (NS-3) & 122,171 & 28.3\\
\dataset{uav\_ew\_bench\_2026} & simulation (physics-informed) & 93,600 & 21.7\\
\dataset{datamut\_sim} & simulation (event) & 15,392 & 3.6\\
\dataset{uav\_attack\_whelan} & real testbed & 610 & 0.1\\
\dataset{uav\_cas} & simulation (digital twin) & --- & ---\\
\midrule
\textbf{real testbed} & captured on hardware & \textbf{200,610} & \textbf{46.5}\\
\textbf{simulation} & generated & \textbf{231,163} & \textbf{53.5}\\
\bottomrule
\end{tabularx}

\end{table}

% =====================================================================
\section{Discussion}
\label{sec:discussion}
% =====================================================================
\emph{What the composition buys an operator.} It turns a network-security knob
into a navigation-safety number. A fleet operator setting a detector's
false-positive budget currently reasons in packet statistics---``how much benign
traffic can I afford to flag?''---with no way to connect that choice to
airworthiness. Theorem~\ref{thm:main} supplies the missing translation: each
detector setting maps to a certified mission floor, so the operator can instead
ask ``what detector tightness do I need to certify my corridor?'' and read the
answer off Fig.~\ref{fig:sensitivity}.

\emph{Why the interface, not the theorem, is decisive.} A reader might expect
the certificate to be the hard part. Our results say otherwise: the tube is a
short consequence of a classical inequality and the Lipschitz constant is
measured, not assumed. What determines whether a useful detector setting is
certifiable is $\gamma$, the metres-per-second of position error a second of
staleness buys the attacker. The order-of-magnitude gap between the kinematic
bound and the measured rate is the difference between a vacuous certified window
and one that contains the operating point. This relocates the research agenda:
tightening cross-layer guarantees is primarily a \emph{measurement} problem
about how attacks propagate into state error---which is precisely why the
benchmark and the theorem belong in one paper.

\emph{Two methodological findings we surface rather than bury.} Building the
benchmark forced two corrections that are instructive for anyone assembling a
cross-layer corpus. First, our initial operating curve was contaminated by an
accumulation bug in the sweep harness, in which each operating point
inadvertently averaged all previous points; the tell was recall \emph{rising}
with the threshold in one regime, which is impossible for a threshold detector.
The clean multi-seed campaign in Fig.~\ref{fig:opcurve} replaces it. Second, the
bus dataset's attack labels are not machine-readable per scenario, so rather
than guess we derived them from frame statistics and flagged the derivation for
confirmation. Both episodes argue \emph{for} the provenance discipline the
schema enforces: a benchmark that records how every number was obtained lets
errors be caught rather than propagated.

\emph{Generality.} Although we instantiate on a time-delay forwarding attack and
a GNSS navigation model, the chain is agnostic to both ends. Any detector with a
monotone operating point and any perturbation-to-state mapping slot into
Lemma~\ref{lem:residual} and \eqref{eq:gamma}; the theorem and the
four-certificate cover are unchanged. Extending the interface catalogue to
jamming-intensity-to-error and to bus-level attacks is direct future work.

% =====================================================================
\section{Limitations and Future Work}
\label{sec:limitations}
% =====================================================================
We state limitations plainly; none undercuts the contributions, and each defines
a concrete extension.

\emph{The delay-to-position rate is a three-flight hover calibration.} The
decisive quantity $\gamma$ is measured on three real flights, all in a low-speed
hover/loiter regime; at cruise speed the position error per second of staleness
could be larger. This is the single result that could change the paper's
framing: if a cruise-inclusive calibration measured
$\gamma>\gamma_{\mathrm{req}}=7.3$\,m/s at a $10$\,m corridor---more than a
five-fold increase over what we observe---the operating point would leave the
certified window and the honest framing would revert to the narrow
``certify-small, detect-large'' regime that the conservative kinematic bound
already guarantees. A larger, cruise-inclusive flight campaign is the
highest-value follow-up.

\emph{Conservatism of the deterministic core.} The Gr\"onwall bound grows as
$e^{LT}$, so for long horizons or higher Lipschitz constants the certified floor
falls below what is empirically observed. This is by design---it is why the
architecture is a four-certificate cover---but fusing the deterministic,
smoothing, and PAC-Bayes floors into a single tighter envelope is open.

\emph{Evidence-class imbalance.} Of the six datasets only one couples a
network/RF event to a measured flight on one airframe, so it is the sole source
of \code{measured\_same\_platform} pairings; the others support physics-modelled
or synthetically-aligned couplings. This honestly reflects what public data
exists, and is exactly why the honesty flag is mandatory rather than optional.

\emph{Calibration constants that remain data-gated.} The unit bridge mapping a
raw GNSS position error into receiver-level feature space (for the
RF-manipulation classes, as opposed to the time-delay class the theorem targets)
is calibrated on a synthetic signal corpus pending access to real in-phase and
quadrature recordings; and the PAC-Bayes bound \eqref{eq:mcallester} is wired to
the trained model's measured empirical risk but awaits its numeric KL
divergence. Neither gates the headline result; both are tracked explicitly so
that no synthetic constant is reported as real.

\emph{Analytical navigation anchor and detector scope.} The benchmark's
completion outcomes come from a physics-informed analytical Monte-Carlo backend
calibrated to full-fidelity anchors; results should be described as a calibrated
analytical benchmark, and we do so. We also instantiate the harness on one
forwarding detector whose operating point is a single clean threshold, chosen
because it makes the cross-layer sweep unambiguous; extending to
multi-dimensional operating points requires a higher-dimensional sweep, not a
schema change.

\emph{Threats to validity.} \emph{Construct:} does a recorded pairing represent
the causal link it claims? The honesty flag bounds this risk, and a reviewer can
restrict any analysis to the measured class. \emph{Internal:} are the numbers
correct? Two of our own measurement errors were caught precisely because the
provenance discipline made them visible, and every figure is regenerated from
committed data. \emph{External:} four of six sources are simulation or
calibrated twin, so absolute detector-performance numbers are corpus-specific;
what transfers is the \emph{structure}---the knee-and-floor operating curve and
the order-of-magnitude gap between kinematic and empirical staleness---which is a
property of store-and-forward scheduling and estimator behaviour, not of any one
simulator.

% =====================================================================
\section{Conclusion}
\label{sec:conclusion}
% =====================================================================
UAV security has been studied as two separate problems---attacks on the wire and
attacks on the flight---because nothing let it be studied as one. We built the
benchmark that makes the coupling representable and then proved, on exactly its
objects, the guarantee that the coupling enables. The benchmark contributes a
two-layer schema whose unit is the cross-layer pairing, six validated adapters,
and a detector-operating-curve harness; the guarantee contributes a residual
budget that measurement tightens into a \emph{saturating} bound
$\Delta(\theta)\le H\min(\theta,s)$, an interface rate calibrated on real
spoofing and jamming flights, and a composition theorem chaining a detector's
operating point to a certified mission-completion floor.

Three findings survive scrutiny. First, the composed guarantee is the only
configuration that certifies a nonzero mission floor against a detector-evading
attacker, dominating both single-layer baselines at every operating point with
overwhelming statistical significance. Second, the empirically calibrated
interface---an order of magnitude tighter than the kinematic worst case---brings
the detector's real operating point inside the certified window, so the
guarantee is operationally relevant and not merely formal. Third, and most
consequential for the field, the binding constraint on cross-layer certification
is the measurement of how attacks propagate into state error, not the proof
machinery: the interface, not the theorem, decides the outcome. UAV network
security and UAV autonomy security, we conclude, are not two problems but two
ends of one certificate.

% NDSS counts 13 pages of BODY; appendices are excluded. This section is
% reader-orientation for reviewers outside UAV autonomy, so it belongs
% here rather than consuming body budget.
% =====================================================================
% =====================================================================
\section{Ethics Considerations}
\label{sec:ethics}
% =====================================================================
This work involves no human subjects and no personally identifying data. All six
datasets are public research corpora released by their authors for security
research; we redistribute none of the raw third-party data, shipping only
integration scripts and unified labels so that each source is fetched from its
origin under its own licence. No live aircraft, production network, or
third-party system was attacked: every attack in this paper is either replayed
from a published capture or generated in simulation on infrastructure we
control.

The dual-use question deserves a direct answer. Our operating-curve harness
quantifies how an attacker can stay under a detector's threshold, which is
information an adversary could in principle use. We judge the disclosure clearly
net-positive and consistent with standard practice in intrusion-detection
research, for three reasons. First, the evasion strategy is not novel---staying
under a published threshold is the definitional behaviour of an evading
attacker, already assumed in the detector's own analysis. Second, our central
result \emph{bounds} what such an attacker achieves and shows it saturates,
which favours defenders. Third, the practical consequence we identify is a
defensive configuration procedure: read the required detector tightness off
Fig.~\ref{fig:sensitivity} for your corridor margin. We disclose no previously
unknown vulnerability in a deployed system, so no coordinated-disclosure process
was applicable; the underlying detector's authors are aware of and have
commented on this analysis of their operating point.

% =====================================================================
\section{Open Science and Artifact Availability}
\label{sec:openscience}
% =====================================================================
Everything needed to reproduce this paper is released as an open artifact: the
schema, the six ingestion adapters, the operating-curve harness, the
certificate engine, the statistical analysis, and the figure generator. Every
figure and table in this paper is machine-generated from committed result files
rather than transcribed, so a reader can regenerate them from the released data;
each carries a provenance tag (real-corpus, simulation-derived, or stated
modelling parameter) in an artifact ledger. The harness is deterministic per
seed, and headline comparisons are released per seed with their nonparametric
tests. We intend to submit the artifact for evaluation.
\ifanon
The repository is anonymised for review and will be de-anonymised on acceptance.
\else
The artifact is available at \url{https://robustuavs.ai}.
\fi

\appendix

\section{Background and Terminology}
\label{sec:background}
% =====================================================================
The composition proved in this paper spans two research communities that share
very few conventions: one reasons about packets and detector thresholds, the
other about trajectories and certified radii. This section fixes the single
system model both halves of the argument refer to---how an autonomous UAV closes
its navigation loop, which of its interfaces an adversary can reach, and what a
robustness \emph{certificate} does and does not promise---so that the interface
mapping of \S\ref{sec:interface} joins two precisely defined quantities rather
than two informal ones. Acronyms are defined at first use and not tabulated;
Table~\ref{tab:notation} is the single reference for mathematical symbols. A
reader fluent in both UAV autonomy and certified robustness may proceed directly
to \S\ref{sec:problem}.

\subsection{How an autonomous UAV stays on course}
An unmanned aerial vehicle (UAV, colloquially a drone) closes a control loop
around a \emph{state estimate}. An onboard computer fuses fixes from a Global
Navigation Satellite System (GNSS, of which the American GPS is one
constellation) receiver, inertial measurements from accelerometers and
gyroscopes (the IMU), and barometric height into a best guess of where the
aircraft is and how fast it is moving. On the PX4 autopilot family used
throughout this paper the fusion is performed by an extended Kalman filter
(EKF)---a standard recursive estimator that weighs each sensor by its expected
noise~\cite{kong2024px4ekf}. The flight controller steers to keep the estimate
on the planned route.

Two properties of this loop drive everything that follows. First, the loop is
only as good as its \emph{freshest trustworthy input}: if GNSS is denied, the
EKF ``coasts'' on inertial data and its position error grows with time. Second,
the loop is \emph{networked}. Swarm members and the ground control station
exchange corrections and commands over radio links: a command-and-control (C2) link, typically
speaking the Micro Air Vehicle Link (MAVLink) protocol, a multi-hop mesh in which peers relay each other's traffic
during periodic contact windows---a delay-tolerant network (DTN)---and an
intra-vehicle bus---the UAV Controller Area Network (UAVCAN, also DroneCAN)---
connecting the flight controller to motor controllers and sensors. An attacker on the mesh can therefore reach the
navigation loop without ever touching the radio-frequency (RF) front end.

\subsection{The two attacker technologies}
\emph{Electronic warfare (EW)} attacks the sensing physics. A \emph{jammer}
raises the jamming-to-signal power ratio ($J/S$, in decibels) until the
receiver loses satellite lock; a \emph{spoofer}
broadcasts counterfeit satellite signals so the receiver locks onto a false
position~\cite{humphreys2008assessing,kerns2014capture}. \emph{Network attacks}
target the messages: a \emph{time-delay attack}---the class our theorem
covers---has a compromised relay hold packets before forwarding them, so
downstream consumers act on stale state~\cite{lou2019timedelay,datamut2026}.
Neither needs to break cryptography; both degrade the control loop through
legitimate-looking inputs, which is precisely why detection and robustness must
be composed rather than assumed.

\subsection{What a ``certificate'' promises}
A \emph{certified} robustness statement differs from an empirical evaluation:
instead of reporting that attacks in a test set were survived, it proves that
\emph{every} perturbation within an explicitly stated budget is survived. The
price is conservatism---certified bounds are worst-case, so a certified floor
typically sits below what is observed empirically. We keep the two kinds of
number separate and label them throughout: \emph{certified} floors come from
proofs, \emph{empirical} rates from measured flights. We use two flavours of
certificate: deterministic Lipschitz-type bounds (if a model cannot amplify
input changes by more than a factor $L$, a bounded input change yields a bounded
output change) and probabilistic randomized-smoothing bounds (adding calibrated
noise and taking a majority vote yields a radius within which the decision
provably cannot flip)~\cite{cohen2019smoothing}. The regulatory anchor for
``mission survived'' is the airworthiness-security process standard
DO-326A/ED-202A~\cite{do326a}, whose industry-acceptable completion floor of
$0.90$ appears in our figures as a horizontal reference line.


\begin{table}[t]
\centering
\caption{Mathematical notation. Every symbol is defined where first used;
this table is the single reference.}
\label{tab:notation}
\small
\renewcommand{\arraystretch}{1.05}
\begin{tabularx}{\columnwidth}{@{}lX@{}}
\toprule
$W_n$, $W_a$ & network / autonomy attack window\\
$P$ & cross-layer pairing tuple, \eqref{eq:pairing}\\
$x(t)\in\Real^n$ & navigation state (position, velocity, \dots) at time $t$\\
$u(t)$, $\delta u$ & sensor input to the navigation dynamics; its perturbation\\
$F$ & drift (right-hand side) of the navigation dynamics \eqref{eq:ode}\\
$T$ & mission horizon (normalised to $1$ in the instantiation)\\
$m$ & safe-corridor margin: allowed lateral deviation, metres\\
$D_\theta$, $\theta$ ($\varepsilon$) & network detector; its operating point
(tolerated per-hop residual delay, s)\\
$s$, $P_{\!c}$ & contact-window slack ($5$\,s); contact period ($60$\,s)\\
$H$ & number of attacker-controlled (malicious) hops on a route\\
$\Delta(\theta)$ & residual budget: worst undetected cumulative delay, s\\
$\delta(\theta)$ & induced perturbation on the navigation input\\
$\gamma$ & delay-to-position rate: metres of error per second of degradation\\
$v_{\max}$ & kinematic worst-case speed bound ($15$\,m/s)\\
$L$, $\hat L$, $\hat L_{\mathrm{loc}}$ & (estimated / local) Lipschitz constant of $F$\\
$\rho(\theta)$ & trajectory-tube radius $\delta(\theta)e^{LT}$, metres\\
$\theta^{*}(m)$ & loosest $\theta$ whose tube fits margin $m$ (window edge)\\
$f(\cdot)$ & certified MCR floor as a function of the perturbation budget\\
$R$, $\sigma,\alpha,n$ & certified input radius; smoothing noise, confidence, samples\\
$R_{\mathrm{emp}}$, $\mathrm{KL}$ & empirical risk and prior-posterior divergence, \eqref{eq:mcallester}\\
$|S|$ & size of the defender's policy set in the MWU certificate ($4$)\\
\bottomrule
\end{tabularx}
\end{table}



\section{Aligning Two Relative Clocks}
\label{app:alignment}
Because the sources share no global clock, a \code{synthetic\_alignment} pairing
must state exactly how a network window and an autonomy window from different
captures are placed on a common axis, and how much error that placement can
introduce. Left implicit, this is the step at which a cross-layer benchmark
could silently manufacture---or destroy---the very effect it reports.

Let the network source have relative clock $\tau_n\in[0,T_n]$ and the autonomy
source $\tau_a\in[0,T_a]$, each measured from its own capture start. An
alignment is an affine map
\begin{equation}
\label{eq:align}
\phi(\tau_n) \;=\; a\,\tau_n + b, \qquad a=\frac{T_a}{T_n},
\end{equation}
whose scale $a$ normalises the two capture durations and whose offset $b$ is
fixed by \emph{anchoring attack onsets}: we set $b$ so that
$\phi(t_0^{\,n}) = t_0^{\,a}$, where $t_0^{\,n}$ is the first labelled malicious
event in $W_n$ and $t_0^{\,a}$ the onset of degradation in $W_a$. Onset is not
eyeballed: $t_0^{\,a}$ is the first sample at which the autonomy metric exceeds
its pre-attack median by more than $3\hat\sigma$, with $\hat\sigma$ the robust
(median-absolute-deviation) scale of the pre-attack segment. Both anchors are
therefore defined by the source labels and the source's own noise floor, not by
the analyst.

Three error terms bound the residual misalignment $e_{\mathrm{align}}$:
\begin{equation}
\label{eq:align_err}
e_{\mathrm{align}} \;\le\; \underbrace{\tfrac{1}{2}(\Delta t_n+\Delta t_a)}_{\text{quantisation}}
\;+\;\underbrace{|a-1|\,T_a}_{\text{rate mismatch}}
\;+\;\underbrace{e_{\mathrm{onset}}}_{\text{anchor detection}},
\end{equation}
where $\Delta t_n,\Delta t_a$ are the two sources' sampling intervals (a per-flow
network record and a $50$\,Hz telemetry stream give
$\tfrac12(\Delta t_n+\Delta t_a)\le\tfrac12(1.0+0.02)=0.51$\,s), the rate term
vanishes when both captures are replayed at their native rate ($a=1$, the case
for every pairing we release), and $e_{\mathrm{onset}}\le\Delta t_a$ is the
one-sample latency of the $3\hat\sigma$ crossing. For the released
\code{synthetic\_alignment} pairings this gives
$e_{\mathrm{align}}\le0.53$\,s.

This bound is what keeps the alignment from biasing the reported staleness. The
quantity the certificate consumes, $\Delta(\theta)$, is computed \emph{entirely
within the network source} from per-hop residual delays; the alignment
determines only which autonomy window a given network window is paired
\emph{with}, never the numeric value of $\Delta(\theta)$. Misalignment can
therefore attach a pairing to the wrong autonomy interval, but it cannot inflate
or deflate $\Delta(\theta)$ itself. Its effect is confined to the empirical MCR
column, and it is bounded there too: since $e_{\mathrm{align}}\le0.53$\,s is two
orders of magnitude below the $60$\,s contact period that separates distinct
attack windows, no released pairing can be displaced onto a neighbouring window.
Pairings whose alignment error would exceed one tenth of the shorter window's
duration are rejected by the adapter rather than emitted with a caveat.

\section{Proof of Theorem~\ref{thm:main}}
\label{app:proof}
\begin{proof}[Proof of Theorem~\ref{thm:main}]
Fix a mission and let $x_{\mathrm{nom}}(t)$ be the nominal trajectory---the
solution of \eqref{eq:ode} under the attack-free input $u(t)$---and
$x_{\mathrm{att}}(t)$ the trajectory under $u(t)+\delta u(t)$, both started from
$x_{\mathrm{nom}}(0)=x_{\mathrm{att}}(0)$. Write
$e(t)=x_{\mathrm{att}}(t)-x_{\mathrm{nom}}(t)$. From \eqref{eq:ode},
$\dot e(t)=F(x_{\mathrm{att}},u+\delta u)-F(x_{\mathrm{nom}},u)$, so
$L$-Lipschitzness gives
$\Norm{\dot e(t)}\le L\Norm{e(t)}+L\Norm{\delta u(t)}$. Integrating with
$e(0)=0$,
\[
\Norm{e(t)}\;\le\;L\!\int_0^t\!\Norm{e(\tau)}\,d\tau
\;+\;L\!\int_0^t\!\Norm{\delta u(\tau)}\,d\tau .
\]
Applying the integral form of Gr\"onwall's inequality~\cite{gronwall1919} to the
non-decreasing forcing term $g(t)=L\!\int_0^t\!\Norm{\delta u}\,d\tau$ yields
$\Norm{e(t)}\le g(t)e^{Lt}$, and bounding the input by
$\Norm{\delta u(\tau)}\le\delta(\theta)$ gives, at the horizon,
\begin{equation}
\label{eq:tube_derivation}
\Norm{e(T)}\;\le\;\delta(\theta)\big(e^{LT}-1\big)\;\le\;\delta(\theta)e^{LT},
\end{equation}
which is \eqref{eq:tube}; we report the closed-form right-hand bound. Now
suppose $\rho(\theta)\le m$. The nominal trajectory lies in the safe corridor by
construction, and every point of the attacked trajectory is within
$\rho(\theta)\le m$ of it by \eqref{eq:tube_derivation}; hence the attacked
trajectory never exits the corridor over $[0,T]$, discharging the spatial
predicate of \eqref{eq:mcr}. For the temporal predicate, the only mechanism by
which a $D_\theta$-evading time-delay attack displaces arrival is the undetected
forwarding delay it accumulates, which Lemma~\ref{lem:residual} bounds by
$\Delta(\theta)$; hence $t_{\mathrm{arr}}\le T+\Delta(\theta)$, and
$\Delta(\theta)\le\kappa T$ gives $t_{\mathrm{arr}}\le(1+\kappa)T$ as required.
Both predicates holding, the mission completes. Taking the probability of
$\{\rho(\theta)\le m\}\cap\{\Delta(\theta)\le\kappa T\}$ over the mission
distribution---the only source of randomness once $\theta$ is fixed, since
$\delta(\theta)$ and $\Delta(\theta)$ are deterministic bounds---gives
\eqref{eq:floor}. Composing
$\theta\!\mapsto\!\Delta(\theta)$ (Lemma~\ref{lem:residual}),
$\Delta\!\mapsto\!\delta$ (\S\ref{sec:interface}), and $\delta\!\mapsto\!\MCR$
establishes the chain \eqref{eq:chain}.
\end{proof}

\ifanon\else

% =====================================================================
\section*{Acknowledgements}
We thank the authors of the \dataset{DATAMUt} detector for the network-side code
and for clarifying its operating point, which grounds the interface mapping, and
the maintainers of the third-party datasets unified here.
\fi

\bibliographystyle{IEEEtran}
\bibliography{refs}

\end{document}
